\documentclass{article}

% Typeset to match Munchausen (arXiv:2007.14430): its own neurips_2020.sty, its package set and
% its theorem environments, so that a page of this diffs against a page of that.
\PassOptionsToPackage{numbers, compress}{natbib}
\usepackage[colorlinks,bookmarks,bookmarksnumbered,linkcolor=blue!55!black,
            citecolor=blue!55!black,urlcolor=blue!55!black]{hyperref}

\usepackage[final]{neurips_2020}

\usepackage[utf8]{inputenc}
\usepackage[T1]{fontenc}
\usepackage{url}
\usepackage{booktabs}
\usepackage{amsfonts}
\usepackage{amsmath}
\usepackage{amssymb}
\usepackage{amsthm}
\usepackage{mathtools}
\usepackage{bm}
\mathtoolsset{showonlyrefs=true}
\usepackage{nicefrac}
\usepackage{microtype}
\usepackage[dvipsnames]{xcolor}

% The style file has no place for a date; its first-page notice does.
\makeatletter
\renewcommand{\@noticestring}{rlexp26 --- 2026-09-21, seventh revision: the distortion family is named,
given the invariants it has to satisfy, and corrected --- the head fits expectiles, not quantiles.
$\srankb$ is written $\srank$ outside the sections where the family is the subject.}
\makeatother

\newcommand{\states}{\mathcal{S}}
\newcommand{\actions}{\mathcal{A}}
\newcommand{\E}{\mathbb{E}}
\newcommand{\h}{\mathcal{H}}
\DeclareMathOperator*{\argmax}{arg\,max}
\DeclareMathOperator*{\argmin}{arg\,min}
\DeclareMathOperator{\mmax}{mellowmax}
\DeclareMathOperator{\sm}{sm}
\DeclareMathOperator{\lse}{lse}
\DeclareMathOperator{\Var}{Var}
\DeclareMathOperator{\Cov}{Cov}
\newcommand{\env}{\mathrm{env}}
\newcommand{\tot}{\mathrm{tot}}
\newcommand{\reg}{\mathrm{reg}}
\newcommand{\gapop}{\delta}
\newcommand{\ta}{\tau_{\mathrm{act}}}
\newcommand{\gap}{\Delta}
\newcommand{\gapbar}{\overline{\Delta}}
\newcommand{\dbar}{\bar{d}}
\newcommand{\sret}{\sigma_{\mathrm{ret}}}
\newcommand{\srank}{\sigma}
\newcommand{\srankb}{\sigma^{\beta}}
\newcommand{\snorm}{\sigma_{G}}
\newcommand{\snoise}{\sigma_{\varepsilon}}
\newcommand{\ip}[2]{\langle #1, #2 \rangle}
\newcommand{\yes}{\ensuremath{\checkmark}}
\newcommand{\no}{\ensuremath{\times}}
\newcommand{\meh}{\ensuremath{\sim}}

\newtheorem{proposition}{Proposition}
\newtheorem{remark}{Remark}

\title{The behaviour policy: what the cases require,\\and the rule that meets them}
\author{rlexp26}

\begin{document}
\maketitle


\begin{abstract}
One object is being chosen: the map from the critic's action values at a state to the
distribution the actor samples from. Today that map is
$\sm(\tilde{q}/\texttt{ACT\_TEMPERATURE})$, and the constant in it carries a value scale, does
not transfer between games, and has no principled setting.

\S\ref{sec:bg} is the setting and every definition used later, including what a family of
distortions has to satisfy and --- \S\ref{sec:family} --- which one this critic turns out to
emit. \S\ref{sec:obj} states what would count as a replacement, and \S\ref{sec:rule} states the
one proposed: $\tau(s) = \min(\srank(s)/g,\; \gapbar(s)/k_0)$, where $\gapbar$ is the mean action
gap at the state and $\srank$ is the spread of the action \emph{ordering} across the critic's own
return distribution --- two reductions over an array \texttt{\_act} already holds, on any critic
that emits a distribution. \S\ref{sec:cases} is the body of the document: ten cases --- seven
constructed, three observed --- each of which imposes one constraint, and a table of which
candidate rules fail which. The cases retire, in order: a fixed sharpness constant, a fixed
absolute regret budget, $\epsilon$-greedy, an entropy target, a calibration read off the top-two
action gap, and a fixed \emph{fractional} regret target.

\textbf{Every number in this document is in \S\ref{sec:meas} and nowhere else.} The sections
before it argue over families of rules and say which properties decide between them; that section
is where a property is given a value, a candidate is given a score, and the family is closed down
to one member. A claim stated qualitatively here is a claim discharged there.

Three of those measurements decide the shape of the rule rather than only its constants.
$\srank$ is chosen over every epistemic denominator because none of the latter transfers, and the
failure is located in the parameter draw rather than in the statistic: \texttt{NoisyLinear}'s
noise is shared across the actor batch and redrawn every step, so it is dithering and not
posterior sampling (Remark~\ref{rem:peraction}). \emph{Nothing computable for free anneals},
which \S\ref{sec:c11} turns into a case that every free rule fails and \S\ref{sec:readmit} is the
price of answering. And coupling $\tau(s)$ into the loss is the better-posed design, on four
counts and against no surviving objection, yet is not the thing to run next --- coupled,
$g \to \infty$ no longer recovers the best arm this project has run, so the family stops
containing it (\S\ref{sec:decouple}). \S\ref{sec:related} is ten papers, of which one moves the
design: Fox (2019) derives \eqref{eq:rule}'s shape in closed form with the signal-to-noise ratio
\emph{squared}, and that exponent is the sharpest open question here.
\end{abstract}

\section{Background and notation}
\label{sec:bg}

Nothing in this section is new and nothing in it is a decision. It is the setting, the four
properties of the agent that the rest of the document uses, every symbol that appears later, and
--- \S\ref{sec:kept} --- two results about M-VI's fixed point that everything below takes as
given.

\paragraph{The setting.} A discounted MDP with a finite action set $\actions$ and
$\gamma = 0.997$, solved by a value-based agent: a critic is fitted by temporal-difference
regression on transitions drawn from a replay buffer, and an actor fills that buffer by playing
the environment. The critic emits numbers, the actor needs a distribution to sample from, and the
map between the two is a free choice. That map is this document's whole subject. The agent is BTR
(\texttt{arXiv:2411.03820}): a dueling IQN head on an Impala trunk, Munchausen, prioritised
replay, and \texttt{NoisyLinear} layers in the head. Four of its properties are used below and
each is stated here.

\paragraph{The critic is distributional.} It emits, per action, a representation of the
distribution $Z(s,a)$ of the discounted return and not only its mean. Two instances matter. IQN
samples quantile levels and returns the return value at each --- $z_u(s,a)$,
\texttt{IQN\_ACT\_SAMPLES} levels per call in \texttt{\_act}; IQN's own name for the level is $\tau$, which is written $u$
throughout so that $\tau$ can be the scale this document chooses. A categorical head instead
returns probabilities $p_i(s,a)$ on a fixed grid of atoms $z_i$. In both,
$\tilde{q}(s,a) = \E Z(s,a)$ is the action value the policy is built from: the mean over sampled
quantiles in the first case, $\sum_i p_i z_i$ in the second. The plan of record puts a categorical
head ahead of this design, so every quantity below is defined as a functional of $Z$ and then
instanced in both. The head is also \emph{dueling}: a state-value stream and an
action-advantage stream, summed. That matters exactly once, at \S\ref{sec:c10}, where it is
what puts risk the action does not control into the value stream and out of the ordering.

\paragraph{Gaps, and the regret of a policy.} $\gap_a(s) = \max_b \tilde{q}(s,b) -
\tilde{q}(s,a) \ge 0$ are the action gaps and $\gapbar(s) = \frac{1}{|\actions|}\sum_a \gap_a(s)$
their mean, which is also the regret of acting uniformly at $s$ and is the \texttt{uniform} row
of \texttt{scripts/probe\_checkpoint.py}. For any policy,
\[
  \rho(\pi;s) \;=\; \max_a \tilde{q}(s,a) - \ip{\pi}{\tilde{q}}(s) \;\ge\; 0
\]
is its \emph{exploration regret} at $s$: what one decision at $s$ costs, in the critic's own
units and on the critic's own account. $Q^\star(s) = \max_a \tilde{q}(s,a)$ is the greedy value
it is measured against, and a regret quoted below as a percentage is $\rho$ over $Q^\star$.
$\h(\pi) = -\ip{\pi}{\ln\pi}$, and $\sm$ is the softmax, spelled as Munchausen spells it.

\paragraph{Distortions, and what a family has to satisfy.} A \emph{distortion} $\beta$ maps a
return distribution to a scalar \emph{on the return axis}: the mean, the median, a quantile, an
expectile, a tail mean. Write $Q_\beta(s,a)$ for the value of $a$ read through $\beta$, and
$A_\beta(s,a) = Q_\beta(s,a) - \frac{1}{|\actions|}\sum_b Q_\beta(s,b)$ for the advantage centred
over actions. A \emph{family} is a set of distortions, written $\beta$ where naming it matters.
Two properties are demanded of every member, and neither is decoration:

\begin{itemize}
  \item \textbf{Translation equivariance}, $\beta(Z + \delta) = \beta(Z) + \delta$. This is what
  delivers the zero set below, and it is the property that fails first. \S\ref{sec:meas-family}
  exhibits two natural readings of a categorical head that lack it --- one of which reports a
  spread that \emph{grows} the better the ordering is resolved, which is the sign error this
  document could least afford.
  \item \textbf{Positive homogeneity}, $\beta(cZ) = c\,\beta(Z)$ for $c > 0$, so that the spread
  below carries the value unit and the gain dividing it stays dimensionless --- which is
  \S\ref{sec:project}'s standing constraint, discharged at the level of the family rather than
  per rule.
\end{itemize}

Quantiles, expectiles, the mean, and a tail mean taken at the distribution's \emph{own} quantile
satisfy both. A tail mean at a threshold fixed on the support does not, and neither does anything
read off atom \emph{probabilities} rather than atom values --- those are not on the return axis
at all.

\paragraph{Two spreads, and what they mean.} Two means over actions:
\[
  \sret(s) = \frac{1}{|\actions|}\sum_a \mathrm{sd}\,Z(s,a),
  \qquad
  \srankb(s) = \frac{1}{|\actions|}\sum_a \mathrm{sd}_\beta\, A_\beta(s,a).
\]
$\sret$ is how wide the returns at $s$ are and needs no family. $\srankb$ is how much the
\emph{ordering} of the actions depends on which part of the return distribution is read: large
when the ordering at the optimistic end disagrees with the ordering at the pessimistic end. It is
a per-state scale, not a per-action bonus --- a mean over actions --- and it is in the units of
value.

\emph{Notation.} The two are contrasted here and nowhere else; past this subsection $\sret$
appears only in \S\ref{sec:c10} and \S\ref{sec:denoms}, where it is the candidate being rejected.
\textbf{The superscript is carried only where the choice of family is the subject --- which is
this subsection, \S\ref{sec:meas-family} and \S\ref{sec:owed} --- and $\srankb$ is written
$\srank$ everywhere else.} A bare $\srank$ therefore always means the ordering spread under
whichever family the head is fitting, $\snorm$ is the dispersion of returns pooled over the
buffer, and $\snoise$ is \texttt{NoisyLinear}'s learned scale; the three are distinct objects and
nothing below shortens any of the other two.

\emph{Its zero set is worth stating exactly, because \S\ref{sec:c10} turns on it.}
$\srank(s) = 0$ iff $A_\beta(s,a)$ is constant in $\beta$ for every $a$, i.e.\ iff
$Q_\beta(s,a) = c_a + m(\beta)$: every action's return distribution is a \emph{translate} of one
common shape. \emph{Translation equivariance is exactly what delivers this} --- under it a
translate moves every reading by the same $\delta$, which the centring then removes, so the
cancellation is an identity rather than an approximation; a family without it reports a spread
where there is nothing to report. That is not an ordering condition and stochastic dominance is neither sufficient
for it nor implied by it --- an action that dominates the rest at every risk level still moves
$\srank$ if its distribution has a different \emph{shape}. So $\srank$ reads departure from a
location family, of which an unresolved ordering is one case and heteroscedasticity across
actions is another.

\paragraph{The two uncertainties, and the estimator that separates them.} A spread of the
critic's output mixes what more data would remove with what it would not. Write
$y_i(\theta)$ for the network's output at quantile index $i$ under parameters $\theta$; then
$\sigma^2_{\mathrm{epi}} = \E_i \Var_\theta(y_i)$ and
$\sigma^2_{\mathrm{ale}} = \Var_i(\E_\theta y_i)$ partition $\Var_{\theta,i}(y_i)$, and both have
\emph{unbiased} estimators from two draws $\theta_A, \theta_B$ of the parameter posterior:
\begin{equation}
  \tilde{\sigma}^2_{\mathrm{epi}} = \tfrac12 \E_i\big[(y_i(\theta_A) - y_i(\theta_B))^2\big],
  \qquad
  \tilde{\sigma}^2_{\mathrm{ale}} = \Cov_i\big(y_i(\theta_A),\, y_i(\theta_B)\big).
  \label{eq:uasplit}
\end{equation}
It is Clements et al.'s and \S\ref{sec:related} is the reading; it is stated here because
\S\ref{sec:c10}, \S\ref{sec:c11} and \S\ref{sec:readmit} all turn on it and because
\S\ref{sec:denoms} reports both halves. Applied to $y_i = z_u$ it splits $\sret$; applied to
the centred advantage it splits $\srank$. One extra pass through the output projections buys
both, and \emph{whether the draws are posterior samples} is the whole question of how much that
is worth --- Remark~\ref{rem:peraction}.

\paragraph{The parameter draw.} \texttt{NoisyLinear} adds $\snoise \odot \varepsilon$ to the
weights and biases of the head's linear layers, with $\varepsilon$ factorised Gaussian noise
redrawn on every call and $\snoise$ learned. It is a large fraction of the network's parameters
and
it is the project's other exploration mechanism: acting greedily on a draw is arm (E) below.
$\dbar(s)$ is the mean over actions of the standard deviation of $\tilde{q}(s,a)$ across draws
of that noise, the quantile sample held fixed --- the previous revision's proposed denominator.
One draw is shared by all $64$ environments and replaced at the next step, so it is dithering and
not posterior sampling; Remark~\ref{rem:peraction} is that reading in full.

\paragraph{Munchausen, and the one coefficient it puts in the loss.} The loss adds
$\alpha\,[\tau_\ell \ln \pi_\theta(a_t|s_t)]_{l_0}^{0}$ to the reward, with
$\pi_\theta = \sm(\tilde{q}/\tau_\ell)$,
$\alpha = \texttt{MUNCHAUSEN\_SCALING\_TERM} = 0.9$,
$\tau_\ell = \texttt{MUNCHAUSEN\_TEMPERATURE} = 0.03$ and
$l_0 = \texttt{MUNCHAUSEN\_CLIPPING\_VAL} = -1$, which is the value every run and every checkpoint
quoted in \S\ref{sec:meas} was produced at. One finished run moved it and it was moved back;
\S\ref{sec:c8-staging} is what that run says, and \S\ref{sec:meas-k} is the one measurement in
this document sensitive to which value is in force. What matters here is that the converged
critic
\emph{is} a policy at the scale $\tau_\ell$ --- its fixed point's gaps are $1/(1-\alpha)$ times
those of the entropy-regularised MDP being solved, so $\sm(q_\infty/\tau_\ell)$ is exactly that
MDP's soft-optimal policy, which is Proposition~\ref{prop:amp} and is proved in
\S\ref{sec:kept} below. $\tau_\ell$ is a coefficient
in the loss and is not what this document chooses; \S\ref{sec:decouple} is why the two are kept
apart for now, and the one place they touch is the ceiling $\tau(s) \le \tau_\ell$ --- a bound
that exists only because they are apart.

\paragraph{Sharpness, and why ``temperature'' is not the word for it.} What a Boltzmann policy
divides by is a scale with the units of value, and every failure below is a failure to get those
units right. The quantity that transfers is its reciprocal in units of the state's own gaps,
\[
  k(s) \;=\; \frac{\gapbar(s)}{\tau(s)},
\]
called the \textbf{sharpness} throughout and quoted in place of a temperature everywhere a
setting is given. It is dimensionless and comparable between games: $k = 0$ is the uniform policy
and $k \to \infty$ is argmax. The status quo is a fixed $\texttt{ACT\_TEMPERATURE}$, so its $k$ is
$\gapbar(s)/\tau$ with $\tau$ held constant --- one constant yielding as many sharpnesses as the
games have gap scales, which is the complaint in one line and which \S\ref{sec:denoms} measures.
$k$ is also a mean, and a softmax does not respond to the mean, which is why it is quoted as a
setting and never used as an estimate of anything.

\paragraph{The constants of the setting.} These are the problem's givens rather than anything
measured, so they are stated here and the theory below is written against them: $\gamma = 0.997$,
so the per-decision horizon is $H = 1/(1-\gamma) = 333$ steps; $n = \texttt{N\_STEP} = 3$;
$|\actions| = 8$ on Phoenix and $18$ on Zaxxon; $64$ environments in the actor; $200$M frames in
a full budget; and Munchausen's $\alpha$, $\tau_\ell$ and $l_0$ as above. Everything else with a
value in it is in \S\ref{sec:meas}. There, every score is quoted as a ratio to BTR's published
number on the same game, and every figure in value units is labelled with its game, because the
two do not transfer into each other.

\subsection{Which family the head is fitting}
\label{sec:family}

\S\ref{sec:bg}'s definitions leave the family open, and it is not a free choice here: the loss
picks it. This is which one is in force under each head, a correction to what the previous five
revisions said about that, and the categorical instance. \emph{Nothing in \S\ref{sec:cases} or
\S\ref{sec:rule} depends on the answer} --- which is why the invariants are stated before the
instance rather than read off it.

\paragraph{Under IQN it is the expectiles, not the quantiles.} The members are the sampled
levels, $Q_u = z_u$, so $\srank$ is the spread of the centred advantage along the level axis and
$\sret$ is the return's own standard deviation; both are reductions over the
$[\texttt{IQN\_ACT\_SAMPLES}, |\actions|]$ array \texttt{\_act} has already formed and is about
to average away. \emph{But $z_u$ is the $u$-expectile.} The quantile Huber is a pinball-weighted
\emph{square} inside its knee and the pinball loss proper outside it, and the minimiser of the
first is the expectile where the minimiser of the second is the quantile; at
\texttt{IQN\_HUBER\_LOSS\_K} the knee sits above the range of $|td|$ a run produces
(\S\ref{sec:meas-family}), so the loss is in its quadratic branch throughout. This document
asserted the quantile reading for five revisions and it was never what the head fitted.

\paragraph{Nothing structural turns on the correction.} Expectiles are translation-equivariant
and positively homogeneous exactly as quantiles are, so \S\ref{sec:bg}'s zero set, the units, and
every case in \S\ref{sec:cases} stand unaltered. What moves is the \emph{scale} of $\srankb$ and
therefore the gain dividing it, which \S\ref{sec:meas-family} prices. One consequence reaches
outside this document: the open \texttt{IQN\_HUBER\_LOSS\_K} item in \texttt{btr.py}'s backlog
proposes lowering the knee to recover the quantile estimand, and doing so would re-scale $g$ ---
a coupling neither that item nor this document had recorded.

\paragraph{Under a categorical head.} $\sret$ is $\sqrt{\sum_i p_i z_i^2 - \tilde{q}^2}$, one
more reduction over the same atoms, and the matching family is the expectiles again, which the
head emits in closed form. The $\beta$-expectile solves
$\beta\,\E[(Z-e)_+] = (1-\beta)\,\E[(e-Z)_+]$, whose two sides differ by a function piecewise
linear in $e$ with knots at the atoms; so with $C_k = \sum_{i\le k} p_i$,
$D_k = \sum_{i\le k} p_i z_i$ and $B_k$, $A_k$ their suffix counterparts, bracketing the sign
change leaves
\begin{equation}
  e_\beta \;=\; \frac{\beta A_k + (1-\beta) D_k}{\beta B_k + (1-\beta) C_k},
  \label{eq:expectile}
\end{equation}
two cumulative sums, one \texttt{searchsorted} and a division --- the cost of a quantile lookup.

\emph{Two claims made in earlier revisions are withdrawn.} The pair ``mean below the median, mean
above it'' is admissible only when each threshold is that \emph{action's own} median, so it is
not ``one dot product with a fixed weight vector''; a threshold fixed on the support is not
translation-equivariant and fails \S\ref{sec:bg}'s first demand. And reading the levels off the
CDF instead --- the direct analogue of $z_u$ --- is admissible but \emph{quantises}: on a fixed
support it is blind below one atom spacing, where \eqref{eq:expectile} is smooth in $p$.
Whether the $\gapbar/\srank$ constant of \S\ref{sec:denoms} survives the change of head is still
\S\ref{sec:owed}'s, but it is now a question about scale rather than about whether the object is
the same one.

\subsection{Two results on M-VI's fixed point}
\label{sec:kept}

Both are short, both were checked numerically, and both bear on constants this document uses.

\begin{proposition}[M-VI's action gap; Munchausen Thm.~2 with the arithmetic redone]
\label{prop:amp}
Let $(q_k,\pi_k)$ be produced by M-VI($\alpha,\tau$) without error, $0 \le \alpha < 1$,
$\lambda = (1-\alpha)\tau$, and $\gapop(q)(s) = \max_a q(s,a) - q(s,\cdot)$. Then
$\lim_k \gapop(q_k) = \frac{1}{1-\alpha}\gapop(q_*^{\lambda})$ and
$\sm(\lim_k q_k/\tau) = \sm(q_*^{\lambda}/\lambda) = \pi_*^{\lambda}$.
\end{proposition}
\begin{proof}
Munchausen's own proof establishes $\lim_k q_k = q_*^\lambda + \alpha\tau\ln\pi_*^\lambda$ and
$\tau\ln\pi_*^\lambda = \frac{1}{1-\alpha}(q_*^\lambda - \lambda\lse(q_*^\lambda/\lambda))$.
Substituting, $\lim_k q_k = (1 + \frac{\alpha}{1-\alpha})q_*^\lambda -
\alpha\tau\lse(q_*^\lambda/\lambda)$, and $1 + \frac{\alpha}{1-\alpha} = \frac{1}{1-\alpha}$. The
log-sum-exp is action-independent, so it drops out of both the gap and the softmax, and
$\frac{1}{1-\alpha}/\tau = 1/\lambda$.
\end{proof}

\begin{remark}[The published constant is $\frac{1+\alpha}{1-\alpha}$ and that is a slip]
\label{rem:erratum}
At $\alpha = 0.9$ the published constant is $19$ against the correct $10$. The error is in the
last line of Munchausen's Appx.~A.4, where $1 + \frac{\alpha}{1-\alpha}$ is collected as
$\frac{1+\alpha}{1-\alpha}$ instead of $\frac{(1-\alpha)+\alpha}{1-\alpha} = \frac{1}{1-\alpha}$;
everything before that line is correct. \texttt{scripts/check\_mvi\_gap.py} iterates M-VI and
entropy-regularised VI to convergence on random finite MDPs and confirms the corrected constant
at several $\alpha$, together with $\sm(q_\infty/\tau) = \pi_*^\lambda$, to the precision
\S\ref{sec:meas-mvi} reports. An independent derivation agrees: decomposing
$Q_\tot = Q_\env + \alpha[\tau\ln\pi]_{l_0}^{0} + C(s)$ gives $\rho_\env = (1-\alpha)\rho_\tot$
exactly, the same factor from the other side.

\emph{Why it is in this document.} It settles what $\tau_\ell$ means, and therefore what the
ceiling of \S\ref{sec:decouple-ledger} is. The network's own output divided by $\tau_\ell$ is the
soft-optimal policy of the MDP being solved; dividing by $(1-\alpha)\tau_\ell$ instead
double-counts the amplification by a factor of $1/(1-\alpha)$, and dividing by
$(1+\alpha)\tau_\ell$, which the published constant implies, is too diffuse by $1+\alpha$. So
$\tau(s) \le \tau_\ell$ is the correct hard bound and not an arbitrary one.
\end{remark}

\paragraph{The distortion, which is the whole of what a coupled $\tau$ changes.} A soft backup at
scale $\lambda$ exceeds a hard one by $\lambda\lse(q/\lambda) - \max_a q$, and that difference is
the entire effect of regularising. It is worth reading rather than carrying, because everything
\S\ref{sec:decouple} weighs is this one number.

\emph{The log-sum-exp is a maximum in disguise.} For any $q$ and any $\lambda > 0$,
\[
  \lambda\lse(q/\lambda) \;=\; \max_{\pi}\big[\ip{\pi}{q} + \lambda\h(\pi)\big],
\]
attained at $\pi = \sm(q/\lambda)$. So the soft backup is not a hard backup with noise added; it
is the best trade the state admits between taking value and keeping entropy, and the softmax is
the policy that strikes it rather than a policy chosen by hand. Subtracting $\max_a q$ from both
sides prices that trade in the two terms this section has already defined,
\[
  \underbrace{\lambda\lse(q/\lambda) - \max_a q}_{\text{what the trade is worth}}
  \;=\; \underbrace{\lambda\h(\pi)}_{\text{credit for the randomness}}
  \;-\; \underbrace{\rho(\pi)}_{\text{regret of collecting it}} .
\]
It is never negative, so \textbf{a Boltzmann policy's entropy credit always covers its regret}.
That is a free bound on the currency of \S\ref{sec:c7}, $\rho(\pi_\tau) \le \tau\h(\pi_\tau)
\le \tau\ln|\actions|$; how slack it is on real states is \S\ref{sec:meas-couple}'s.

Coupling \S\ref{sec:rule}'s rule means putting $\tau(s)$
in place of $\tau_\ell$ in the loss, hence $\lambda(s) = (1-\alpha)\tau(s)$, and then the agent's
objective gains that quantity at every step. \emph{Read on the network's own output it carries no
$\alpha$ at all}: by Proposition~\ref{prop:amp} the gaps of the regularised MDP are $(1-\alpha)$
times the ones the critic emits, so the $(1-\alpha)$ in $\lambda$ cancels inside the exponent,
and the amplification $\frac{1}{1-\alpha}$ cancels the one outside it. What is left is
\begin{equation}
  D(s) \;=\; \tau(s)\,\ln\sum_a e^{-\gap_a(s)/\tau(s)},
  \qquad 0 \;\le\; D(s) \;\le\; \tau(s)\ln|\actions|,
  \label{eq:distortion}
\end{equation}
in the units of $\tilde{q}$ and directly comparable with $\gapbar$ --- and the decomposition above
survives the cancellation unchanged, $D(s) = \tau\h(\pi_\tau) - \rho(\pi_\tau)$ read at the
critic's own scale, which is the form to check an implementation against. It is what
\S\ref{sec:c8} weighs and \S\ref{sec:meas-couple} measures.

\emph{In words.} $D(s)$ is what the state's remaining choice is worth --- the premium a soft
backup pays for not yet having to commit. It is zero where one action is clearly best, because
there is nothing to hedge, and largest, $\tau\ln|\actions|$, where the actions are tied and the
choice is free. Coupled, that premium is added to the return at every step, so the agent is paid,
once per visit, for standing where its own critic has not yet decided. Whether that is a useful
thing to pay for is \S\ref{sec:c8}; what it costs is \S\ref{sec:meas-couple}.

\emph{How it behaves.} Exponentiating \eqref{eq:distortion} names the quantity $D$ is really
about. Write
\[
  m(s) \;=\; e^{D(s)/\tau(s)} \;=\; \sum_a e^{-\gap_a(s)/\tau(s)} \;\in\; [1,\,|\actions|],
  \qquad\text{so}\qquad D(s) \;=\; \tau(s)\ln m(s):
\]
the greedy action contributes exactly $1$ and every other action contributes according to how
far it is behind in units of $\tau$, so \textbf{$m$ is the number of actions still in contention
at that scale} and $D$ is the temperature times its logarithm. The two limits are immediate:
$m \to 1$ where one action is clearly best and $D \to 0$; $m \to |\actions|$ at a full tie and
$D \to \tau\ln|\actions|$.

\emph{The top-two shortcut is a special case and not a licence.} If exactly one runner-up is in
contention then $m \approx 1 + e^{-k_2}$ with $k_2 = \gap_2/\tau$ the sharpness at the top-two
gap, and $D \approx \tau e^{-k_2}$ --- linear in the temperature, exponential in the near-tie.
That form is worth having because it shows the exponential does the deciding, but \emph{it is
only correct when the sum has one live term}, and $D$ is largest exactly where it may not:
\S\ref{sec:meas-couple} scores it where the bonus actually lives rather than at a typical state,
and finds it materially wrong there and worse on the larger action set. \textbf{$m$, not $k_2$,
is the quantity to reason with.} The mean sharpness $k(s)$ is wronger still: it averages over all
actions, including the ones $D$ cannot hear.

Because $D$ is a logarithm of a sum of exponentials, $D/\tau$ ranges over dozens of decades
across states, so a mean of it is a tail statistic rather than a typical value --- which is why
\S\ref{sec:meas-couple} reports $D$-weighted means rather than correlations.

\emph{And $m$ counts actions, which makes $D$ sensitive to how the action set is written down.}
Duplicating every action multiplies $m$ by the number of copies exactly, so
\[
  D \;\longmapsto\; D + \tau(s)\ln(\text{copies}),
\]
with no appeal to approximation. The \emph{policy} is untouched by this --- a softmax spreads the
same total mass over each duplicated value, which is \S\ref{sec:c3} --- so the invariance the
behaviour rule has, the objective a coupled $\tau$ induces does not. \S\ref{sec:c8} is where that
matters, because the added term is proportional to $\tau(s)$ and therefore to whatever the rule
divides by.

\begin{proposition}[Regret is monotone in the scale, and priced by any value function]
\label{prop:rho}
Let $\pi_\tau = \sm(v/\tau)$ on a finite action set for any $v$, and
$\rho_w(\tau) = \max_a w(a) - \ip{\pi_\tau}{w}$ for any $w$. Then
$d\rho_w/d\tau = \Cov_{\pi_\tau}(v,w)/\tau^{2}$. At $w = v$ this is
$\Var_{\pi_\tau}(v)/\tau^{2} \ge 0$.
\end{proposition}
\begin{proof}
With $b = 1/\tau$, $\partial\pi_a/\partial b = \pi_a(v_a - \ip{\pi}{v})$, so
$d\ip{\pi_\tau}{w}/db = \Cov_{\pi_\tau}(v,w)$ and
$d\rho_w/d\tau = -(d\ip{\pi_\tau}{w}/db)(db/d\tau) = \Cov_{\pi_\tau}(v,w)/\tau^{2}$.
\end{proof}

Verified against a finite difference at several $\tau$ on unrelated random $v, w$. The $w = v$
case is the heat capacity of a Boltzmann distribution, $C = \Var(E)/T^{2}$, with $-v$ for energy.
It is kept because it is what makes the $\rho$ series of \S\ref{sec:meas-rho} interpretable:
$\rho$ increases in $\tau$, so a measured $\rho$ inverts to a sharpness and the four checkpoints'
tables can be read against each other. It is no longer inverted numerically by anything --- that
was the root-find \S\ref{sec:c9} and \S\ref{sec:rule} removed.


\section{Problem statement and criteria for adequacy}
\label{sec:obj}

\paragraph{The object.} A map $\tilde{q}(s,\cdot) \mapsto \pi(\cdot|s)$, evaluated in
\texttt{\_act} once per environment step per environment, and read by nothing else. Everything
downstream --- the loss, the target, the replay buffer, the priorities, the evaluator --- is
untouched by this document.

\paragraph{The criterion, in one sentence.} A rule is acceptable if its constants can be set once
and left alone across games, across a run's own movement in \texttt{train/q}, and
across the two failure modes this project has observed --- and if it costs nothing to run and
nothing to revert.

That is a lower bar than ``explores optimally'' and it is set there deliberately, because
\S\ref{sec:meas-return} shows the measurable prize is small: four full-budget arms spanning a
decade of sharpness converge to within a small fraction of each other late in a run, which is
less than one seed is worth. What is \emph{not} small is the early term, where the same four arms
span several-fold (\S\ref{sec:meas-return}), and every one of the cases below bears on it. So
the design goal is not a
better exploration--exploitation trade-off. It is a behaviour policy whose failure modes are
bounded by construction rather than by a constant somebody chose, at no cost, that reduces to the
best-measured arm when the value function is informative.

\paragraph{What the criterion excludes at this price.} A second pass through
\texttt{value\_linear0}/\texttt{advantage\_linear0} (most of the network's parameters), a
re-established full-budget baseline, and anything whose constants are in value units.
\S\ref{sec:dropped} lists what that removes and why each removal is safe, and the rule of
\S\ref{sec:rule} comes in under it with nothing to spare: two reductions and a division.

\emph{It does not exclude a head or a servo outright, and the previous revision was wrong to
read it that way.} What it excludes is paying for one without a case that demands it.
\S\ref{sec:c10} and \S\ref{sec:c11} are two such cases, both failed by every rule that runs for
free, and \S\ref{sec:readmit} is what they buy.

\subsection{Constraints external to the case analysis}
\label{sec:project}

They are properties of this project rather than of the problem, and they are stated first because
each one disqualifies candidates before any case is heard.

\emph{The control path is the asset.} There is one full-budget reference per game and it was
expensive. A change that perturbs the loss forfeits the comparison until it is re-run. A change
confined to \texttt{\_act} does not, and can be reverted by deleting it. This is the asset
\S\ref{sec:decouple} weighs a coupled loss against, and the only reason it defers one.

\emph{No constant may carry a value scale.} \texttt{train/q} moves by more than two orders over
a run, \texttt{train/action\_gap} by one, and $\mathrm{gap}/q$ is measured to \emph{fall} rather
than hold; across games the two differ by a further large factor (\S\ref{sec:denoms},
\S\ref{sec:meas-gap}). A constant in value units is therefore wrong at two ends of one run and at
both games, and no
amount of tuning fixes it. This is the whole complaint against \texttt{ACT\_TEMPERATURE} and any
replacement that repeats it is not a replacement. There are two ways to satisfy it: quote every
constant as a ratio of two quantities the network emits, which is what \eqref{eq:rule} does, or
divide the value unit out globally with a return normaliser, which is what the primed
--- return-normalised --- candidates of \S\ref{sec:cases} do and \S\ref{sec:norm} prices.

\emph{Nothing may depend on how the critic is implemented.} The plan of record puts a categorical
head ahead of this design, and \S\ref{sec:fixdraw} contemplates deleting \texttt{NoisyLinear}
altogether, which would return a large fraction of the parameters. A rule that assumes sampled
quantile
levels, or assumes a parameter draw is present to do the exploring, is a rule that has to be
redesigned when either lands --- and worse, a rule whose constants would have to be re-measured
then, which is the same failure as a constant in value units wearing a different coat. So every
quantity used below is a functional of the return distribution the critic emits
(\S\ref{sec:bg}), stated once and instanced twice; and every rule is required to remain a
well-defined policy when the parameter draw is removed. Stating it as a constraint here is what
removes the previous revision's separate section arguing, after the fact, that its rule survived
a change of head.

\section{The proposed rule}
\label{sec:rule}

\begin{equation}
  \boxed{\;
  \tau(s) \;=\; \min\!\left(\frac{\srank(s)}{g},\; \frac{\gapbar(s)}{k_0}\right),
  \qquad
  \pi(\cdot\,|\,s) \;=\; \sm\!\left(\frac{\tilde{q}(s,\cdot)}{\tau(s)}\right)
  \;}
  \label{eq:rule}
\end{equation}

\paragraph{In words.} \emph{Act at the width of the state's own ordering, divided by a gain, and
never more diffusely than a floor.} Both constants are dimensionless: $g$ divides one spread by
another, and $k_0$ is a sharpness. They are calibrated at \S\ref{sec:meas-g} and
\S\ref{sec:meas-k}, which is where they are given values. In the sharpness of \S\ref{sec:bg} the
same rule reads
\begin{equation}
  k(s) \;=\; \max\!\left(g\,\frac{\gapbar(s)}{\srank(s)},\; k_0\right),
  \label{eq:rule-k}
\end{equation}
which is the form to compare against the arms that have run, every one of which is a $k$.

\paragraph{Where each piece comes from.} Every line below is a case in \S\ref{sec:cases}; the
case is the argument and this is only the index into it.
\begin{itemize}
  \item \textbf{A Boltzmann policy} rather than a rate: \S\ref{sec:c7}, and \S\ref{sec:c2} in its
  extreme form. Pricing actions by value is how a given off-greedy rate is bought cheaply, and
  $\epsilon$-greedy pays an order of magnitude more for the same rate on the measured
  checkpoints.
  \item \textbf{$\srank$ in the denominator}: \S\ref{sec:c5} for what it has to carry,
  \S\ref{sec:c10} for why not $\sret$, \S\ref{sec:denoms} for the fact that it transfers and
  $\dbar$ does not. It is the only quantity measured here that carries the state's value scale
  \emph{and} how far the ordering is from resolved, and carries them in one number because it is
  a spread of the critic's own output --- which is also why \S\ref{sec:c1} and \S\ref{sec:c6}
  come out exact rather than conditional.
  \item \textbf{Means over actions, not top-two gaps}: \S\ref{sec:c3}. Both $\gapbar$ and
  $\srank$ are duplication-invariant; the top-two gap is not, and \S\ref{sec:denoms}'s table has
  it far apart between this project's two games largely for that reason.
  \item \textbf{$\gapbar$ only as a floor, never as a target}: \S\ref{sec:c2} and \S\ref{sec:c4}
  say a target proportional to $\gapbar$ spends most where the options are worst. Its job is
  narrower than in the previous revision --- it is a guard against $\srank$ itself failing, not
  the rule's answer to \S\ref{sec:c6} --- and \S\ref{sec:meas-k} places it where fewer than one
  state in twenty rests on it.
  \item \textbf{A $\min$ of two bounds rather than a blend}: the two are different kinds of
  statement and there is no exchange rate between them. Both are upper bounds on how diffuse the
  policy may be, so both are enforced.
  \item \textbf{A closed form rather than a root-find}: $\srank$ and $\gapbar$ are available
  exactly from an array already in registers, so there is no bracket, no bisection count and no
  tolerance. The one line this deserves is that the constraint form has a degeneracy where the
  closed form has none, which is \S\ref{sec:c9}. Where \S\ref{sec:readmit} later takes a
  parameter, it is for a quantity the critic does not emit at all; nothing there is about
  approximating this one.
\end{itemize}

\paragraph{Cost.} The trunk runs once and is unchanged, and no second draw is taken. $\srank$ and
$\gapbar$ are two reductions over the $[\texttt{IQN\_ACT\_SAMPLES}, |\actions|]$ array
\texttt{\_act} already holds, then a division, a $\min$ and the softmax that is there anyway; a
categorical head pays two extra dot products instead (\S\ref{sec:bg}).
\textbf{Zero parameters, zero optimiser state, zero heads, zero target networks, no extra pass,
no root-find.} In the learner: nothing at all while the two are decoupled, and two more
reductions over arrays \texttt{\_train\_step} already holds if they are not
(\S\ref{sec:couple-fixedpoint}). This is strictly cheaper
than the previous revision's rule, which wanted a second pass through the output projections to
form $\dbar$.

\paragraph{Which way it fails.} Under IQN, $\srank$ is built out of the head's
$u$-modulation, and the LayerNorm-plus-weight-decay run is measured to anneal
\texttt{quantile\_embedding}. If that modulation decays, $\srank$ decays with it and the rule
sharpens to argmax --- which is the best-measured arm, where the status quo's failure direction
is the uniform policy, which \S\ref{sec:meas-return} measures as the larger loss. That asymmetry
is not an accident of
this denominator but it is worth naming, and $\srank$ is logged so that the decay is visible
rather than silent.

\paragraph{It contains the best-measured arm.} As $g \to \infty$, $\tau(s) \to 0$ and
\eqref{eq:rule} is argmax acting with the parameter draw --- arm (E), which
\S\ref{sec:meas-return} scores above every soft arm over the second half of a full budget. The
family's sharp limit is the best thing this project has run, so the question a run answers is
whether a finite $g$ buys anything on top of it, not whether the apparatus beats the baseline.
That is the only staging under which a negative result costs one run and no baseline.

\begin{remark}[The parameter draw is not posterior sampling as it is wired]
\label{rem:peraction}
The previous revision claimed the codebase already performs approximate Thompson sampling, since
drawing the head noise and acting greedily on the draw is posterior sampling. The claim is
withdrawn. \texttt{NoisyLinear} draws $\varepsilon_{\mathrm{in}}$ of shape
$(\texttt{in\_features},)$ and $\varepsilon_{\mathrm{out}}$ of shape
$(\texttt{out\_features},)$, with no batch axis, on every \texttt{\_\_call\_\_}: one perturbation
is shared by the whole actor batch, and it is replaced at the next environment step. Posterior
sampling requires a draw that is independent across actors and \emph{held} while its consequences
play out; this is neither, so it is dithering with a value-weighted shape, and the deep
exploration the argument assumed is not present to be built on. What survives is the magnitude:
\S\ref{sec:meas-rho} finds the draw accounting for most of the best soft arm's total exploration
regret, so \eqref{eq:rule} is a second, undirected layer on top of a first whose character has
been mis-described. Two nearly free changes would make the description true, and
they are \S\ref{sec:fixdraw}'s.
\end{remark}

\section{Case analysis}
\label{sec:cases}

Ten cases: seven constructed, three this project has observed. Each is stated, then what it
demands, then which candidates fail it. \eqref{eq:rule} is on the list as (F), already stated, so
that a case can be read as a verdict rather than as a promise.

\paragraph{The candidates.} Each is a map from $\tilde{q}(s,\cdot)$ to a distribution over
actions.

\begin{description}
  \item[(A) A fixed sharpness constant.] $\pi = \sm(\tilde{q}(s,\cdot)/\tau)$ with $\tau$ a
  number in the source. The status quo, at $\tau = \texttt{ACT\_TEMPERATURE}$.
  \item[(A$'$) The same, in normalised units.] $\tau = c\,\snorm$, where $\snorm$ is a
  running estimate of the dispersion of returns in the buffer --- PopArt's scale --- and $c$ is
  dimensionless. Everything (A) does, with the value unit divided out once per game instead of
  never.
  \item[(B) $\epsilon$-greedy.] The greedy action with probability $1-\epsilon$, uniform
  otherwise. BTR's own, at a small floor $\epsilon$ after a geometric decay.
  \item[(C) A fixed absolute regret budget.] The most diffuse $\pi$ with
  $\rho(\pi;s) \le \rho^\star$, for $\rho^\star$ a constant in value units.
  \item[(C$'$) The same, normalised.] $\rho^\star = c\,\snorm$.
  \item[(D) A fixed fractional regret budget.] The most diffuse $\pi$ with
  $\rho(\pi;s) = c\,\gapbar(s)$: spend a fixed fraction of the regret that is available at the
  state. The previous revision's recommendation.
  \item[(E) Argmax.] $\argmax_a \tilde{q}(s,a)$ evaluated on a draw of the parameter noise, which
  is then the only exploration in the loop. The best-measured arm.
  \item[(F) \eqref{eq:rule}.] $\tau(s) = \min(\srank(s)/g,\ \gapbar(s)/k_0)$.
  \item[(I) A contracting scale.] $\tau(s) = \hat{s}(s)/g$ for any per-state signal $\hat{s}$
  that shrinks as a state is visited: the epistemic spread of a posterior over the action values,
  the disagreement of an ensemble, or a novelty score carried into value units. (I) is a family
  and not a proposal --- it is carried so that the last two cases have something that passes them
  --- and what an instance costs is \S\ref{sec:readmit}.
\end{description}

An entropy target \textbf{(G)}, $\h(\pi)(s) = \bar{\h}$, which is SAC's constraint, is carried
until \S\ref{sec:c3} decides it. \emph{A primed candidate inherits its sibling's verdict except
where a case says otherwise}, which happens three times: \S\ref{sec:c1}, where the normaliser is
the whole point, and \S\ref{sec:c5} and \S\ref{sec:c6}, where it is measurably not enough.

For (C), (C$'$) and (D), note once that a maximum-entropy policy subject to a linear constraint
$\E_\pi[\gap] = \rho^\star$ \emph{is} a Boltzmann distribution, with the scale as the
constraint's Lagrange multiplier. So they are not alternatives to a softmax; they are a softmax
whose scale is solved per state rather than chosen. That framing is what makes them comparable to
(A) on the same axis, and it is one more reason that whichever scale wins should reach the loss
as a temperature and not as a solved constraint (\S\ref{sec:decouple}): a dual variable is not a
modelling coefficient.

\subsection{Invariance under reward rescaling}
\label{sec:c1}

\emph{Demand.} The behaviour policy is unchanged. Nothing about the decision problem changed;
only the unit did.

(A) fails outright --- $\sm(100\tilde{q}/\tau)$ is a hundredfold sharper --- and this is the
\texttt{ACT\_TEMPERATURE} complaint in its purest form. (C) fails the same way from the other
side: a fixed regret budget is a hundredth of the available regret after rescaling, so the policy
sharpens. (B) and (D) pass by construction, being defined on rates and ratios respectively.

\emph{This is the case the primed candidates exist for, and they pass it.} $\snorm$ is a
dispersion of returns, so it scales with the reward unit and $\tau = c\,\snorm$ and
$\rho^\star = c\,\snorm$ both scale with it. The pass is exact in the limit and approximate in
practice, since $\snorm$ is an estimate from returns already in the buffer and therefore lags
a change of unit by however long the buffer is --- which is a real effect over a run, where the
scale a game's returns are measured on moves as the agent gets better at it. \S\ref{sec:norm}
prices the normaliser; this case is what buys it.

(F) passes exactly and unconditionally: $\srank$ is a spread of the critic's own outputs, so it
is homogeneous of degree one in the reward unit, and $\tilde{q}/\srank$ is invariant. This is the
one place the change of denominator buys something for nothing --- the previous revision's
$\dbar$ passed here only \emph{conditionally}, on $\snoise$ tracking the value scale, because the
parameter noise is a separate quantity that the reward unit does not reach. \S\ref{sec:denoms}
measured that condition and it does not hold.

(E) inherits the same conditional and keeps it, the draw being the whole of its exploration. Two
things soften that: the failure is toward \emph{sharpness}, which costs return at the margin and
cannot destabilise, where (A)'s and (D)'s failures below are toward diffuseness; and (E) is the
best-measured arm as it stands.

(I) passes when its signal is in value units --- an epistemic spread of the action values is, and
rescales with them --- and fails when it is not. A count, a pseudo-count or a novelty score is a
number about states, not about returns, so an instance built on one needs a carrier to multiply
it into value units, and then it is that carrier and not the signal that has to pass this case.
This is the $\ddagger$ on (I) throughout, and \S\ref{sec:contracting} is where the carrier is
chosen.

\subsection{Invariance under duplicated actions}
\label{sec:c3}

\emph{Demand.} The marginal behaviour over distinct actions is unchanged, and the constant does
not have to be retuned. This is not hypothetical: Zaxxon's action set is five directions crossed
with fire and aliases heavily, which is most of why \S\ref{sec:denoms} finds its top-two gap so
far below Phoenix's --- a factor that is partly the action set and not the task.

Every regret-based rule passes. $\gapbar$ is invariant under duplication --- each value is
counted twice in a mean over twice as many actions --- and a Boltzmann policy's marginal over
distinct actions is the Boltzmann policy over distinct actions at the same scale. So (A), (B),
(C), (D), (E), (F) all pass, and so do the primed variants, $\snorm$ being a property of
returns and not of the action set. (I) passes for the same reason as (F): every statistic in it
is a mean over actions, and whatever a posterior or an ensemble says about an action it says
identically about that action's twin.

What this case decides is what is \emph{not} on the list. \textbf{(G), an entropy target, fails}:
$\h$ ranges over $[0, \ln 2m]$ where the decision problem has $m$ distinct actions, so a target
in nats is a different policy before and after duplication, and SAC's $\bar{\h}$ is exactly such
a target. And \textbf{a calibration read off the top-two gap fails}: a duplicated argmax drives
that gap to zero at unchanged behaviour, so $\tau = \mathrm{gap}_{\mathrm{top2}}/k$ --- the rule
\texttt{docs/plan.md} currently proposes --- collapses. Both are excluded here and do not appear
again.

\subsection{Added dominated actions}
\label{sec:c4}

\emph{Setup.} The same decision problem with $16$ further actions bolted on, all worse than
everything already there: one best action and the rest at a common gap $\delta$.

\emph{Demand, stated carefully, because the obvious statement of it is wrong.} Sixteen unpriced
actions have to be tried before anything can be said about them, and trying them costs
exploration --- that cost is the problem's and every rule pays it. A Boltzmann policy pays it in
the cheapest available order, spending most on whatever it currently rates highest among the
unknowns; that is what value-weighting is for and no case here objects to it. What a rule may not
do is \emph{go on} paying once the answer is in. So the demand is at fixed knowledge: with the
critic having resolved that the added actions are worse, enlarging the action set must not make
the policy more diffuse than it was on the original set.

This is \S\ref{sec:c3}'s matched pair and it is where the two part company. With one best action
and the rest at a common gap $\delta$, $\gapbar = \delta(|\actions|-1)/|\actions|$: that is
$0.500\delta$ at $|\actions| = 2$, $0.875\delta$ at $8$ and $0.944\delta$ at $18$. So
\textbf{(D) fails}. Its target $c\,\gapbar$ inflates by $1.75\times$ from two actions to eight,
and since regret is spent at $\delta$ per unit of off-greedy mass, the off-greedy probability
inflates with it --- permanently, because the added actions stay bad and so $\gapbar$ stays
inflated. The rule explores more because more bad options exist, long after it has established
that they are bad.

That is the same objection the previous revision raised against entropy and against the top-two
gap --- that they are functionals of the \emph{action set} rather than of behaviour --- turned
back on the statistic it then built its own target out of. $\gapbar$'s invariance is exactly
duplication-invariance and no more.

(C) and (C$'$) pass: a budget in value units does not read the action set at all. (B) passes on
the stated demand, its off-greedy rate being $\epsilon$ regardless, though the \emph{cost} of that
rate rises as the added actions are worse and uniformly drawn --- which is \S\ref{sec:c2}, not
this case. (A), (A$'$), (F) and (I) are all mildly exposed: more actions means more terms in the
softmax's tail, so the off-greedy rate rises linearly in the count at fixed scale. That is an
acceptable residue and it is the standard behaviour of a Boltzmann policy, not a defect
introduced by this design. Whether (F) is exposed \emph{more} than (A) depends on what the added
actions do to $\srank$, which is a mean over actions and so does move when the action set does;
confidently bad actions have narrow orderings and pull it down, which sharpens, but this has not
been measured and the verdict is the same either way. (I) moves the same way and for the same
reason: a resolved bad action has little epistemic spread, so the mean falls and the policy
sharpens.

\subsection{Catastrophic actions}
\label{sec:c2}

\emph{Setup.} $|\actions| = 2m$; $m$ actions return $-M$ immediately for large $M$; the other $m$
differ among themselves by a small $\delta$ and are what precise control is about.

\emph{Demand.} The probability of taking a catastrophic action falls as $M$ grows. A rule that
takes it at a rate independent of $M$ is not making a trade-off, it is ignoring one.

$\gapbar \approx M/2$, dominated entirely by the actions to be avoided. Regret can only be spent
in quantity by putting mass on them, so a target of $c\,\gapbar$ forces $p_{\mathrm{bad}}\cdot M
\approx c M/2$, that is \textbf{$p_{\mathrm{bad}} \approx c/2$ regardless of $M$}: (D) fails,
and it fails for a structural reason --- $\gapbar$ is inflated by precisely what makes the state
dangerous, so the rule reads ``much regret is attainable here'' and spends it. (B) fails
identically at $p_{\mathrm{bad}} = \epsilon/2$.

(A) passes, $p_{\mathrm{bad}} \propto e^{-M/\tau} \to 0$; this is value-weighting doing its job
and it is what a Boltzmann policy is for. (C) passes and is the most precisely right of them:
$p_{\mathrm{bad}} \approx \rho^\star/M$, so the mistake gets rarer at exactly the rate that keeps
its cost constant. (E) passes. (F) passes, because its operating branch does not read $\gapbar$
at all: $\tau = \srank/g$, and an action that is confidently catastrophic has a return
distribution that is confidently at $-M$, contributing nothing to the spread of the ordering. So
$p_{\mathrm{bad}} \propto e^{-Mg/\srank} \to 0$, and the floor branch does not bind,
$\gapbar/k_0 \approx M/2k_0$ being enormous. (I) passes by the same two steps: a resolved
catastrophe has no epistemic spread either, and the exponential does the rest.

\emph{This case and \S\ref{sec:c4} together retire (D)}, which is the previous revision's
recommendation, and they retire it on the same mechanism seen twice: a target proportional to the
attainable regret is a target proportional to how bad the worst options are.

\subsection{Two regions at different reward scales}
\label{sec:c5}

\emph{Setup.} Room 1 rewards precise control and the agent has learned it; room 2 rewards
exploration, tolerates mistakes, and its rewards are $1000\times$ larger. One game, one network,
two regions of its state space.

\emph{Demand.} Sharp in room 1, diffuse in room 2, from one setting.

(A) fails hardest, and this is the case the whole per-state idea exists for: one $\tau$ against
gaps that differ by $1000\times$ is near-argmax in one room and near-uniform in the other, and no
value is right for both. That is not a thought experiment on this codebase --- it is the
value collapse \S\ref{sec:meas-collapse} records --- a Phoenix run whose \texttt{train/q} fell
while the policy went near-uniform --- in miniature.

\textbf{(A$'$) fails too, and this is the limit of normalisation.} $\snorm$ is one number for
the whole game, estimated from returns pooled over both rooms, so $\tau = c\,\snorm$ is one
temperature again --- the right one nowhere and wrong by up to $1000\times$ somewhere. A global
normaliser converts a constant that is wrong across games into a constant that is wrong across
states within a game, which is progress and is not the case. (C$'$) inherits it: its budget is
infeasible in the poor room for the same reason (C)'s is.

(C) is sharp in the rich room and infeasible in the poor one ($\rho^\star > \gapbar$ has no
solution, and the previous revision's simulation finds its sharpness near-uniform and
essentially uncorrelated with the state once the gaps contract). (B) and (D) are \emph{equally}
sharp in both rooms, which is not what the case asks for but is bounded: they turn a
$1000\times$ error into a constant-factor one, which is the honest size of what scale-freeness
alone buys.

\textbf{(F) passes, and so does (I)}; the reason is worth isolating because it is the document's
central claim. What the case actually asks is that sharpness track \emph{how much exploring is
worth}, and that is not a property of the value scale at all. $\srank$ carries both: it is a
spread of that room's own values, so it moves with the room's reward scale, \emph{and} it is the
spread of the \emph{ordering}, so it moves with how far the network is from having resolved which
action is best there. Room 1 is learned, so the gaps are large against $\srank$ and the policy is
sharp; room 2 is not, so they are not and it is diffuse. Neither response is tuned. (I) passes
for the second half of that sentence alone, and more directly: an epistemic signal is exactly a
measure of how much is left to learn in the room.

The mechanism is not assumed here, which it was in the previous revision. $\gapbar/\srank$ reads
measurably higher on the argmax-trained net than on the soft-trained one
(\S\ref{sec:denoms}) --- so the net whose ordering is better resolved, and which scores higher, is
the one the rule acts more sharply on, with no constant changed.
That is this case's demand observed between two real checkpoints rather than between two
imagined rooms.

\emph{What $\srank$ does not carry} is which of the two reasons a room is diffuse: a state whose
ordering is unresolved because it is unvisited and one whose ordering is unresolved because the
return is genuinely noisy look the same to it. That is \S\ref{sec:c10}, and it is the price of a
denominator that costs nothing.

\subsection{Uninformative value functions: initialisation and collapse}
\label{sec:c6}

\emph{Setup.} All $\tilde{q}$ are near-equal --- at initialisation, or after a value collapse.

\emph{Demand.} The policy does not become uniform, and in particular the rule does not
\emph{amplify} the contraction: the policy it produces at contracted gaps is the policy it would
produce at the same gaps scaled up.

\emph{Why uniform is the wrong answer here, since it is not obvious.} The objection to make is
that a critic which knows nothing cannot direct anything, so a uniform random walk is the best
available exploration. For a bandit, or a short horizon, or an agent with an intrinsic-motivation
term to direct it instead, that is correct. It is not correct here, for two reasons that are
properties of the problem and not of the rule. First, reward in these games is reached by
committed sequences of actions rather than by single ones, so a policy that redraws uniformly at
every one of $H = 333$ steps reaches it with probability $|\actions|^{-L}$ for a sequence of
length $L$ --- zero in practice for Enduro's ``hold accelerate'', and small enough on Phoenix to
matter. The returns such a policy collects differ only by noise, the critic fits that, the gaps
do not grow, and the loop closes: Munchausen's Appx.~B.2 reports exactly this state on Enduro
($Q$ near zero, a softmax ``very close to uniform''), that the agent gathers no reward, receives
no signal, and never recovers. This project has the mid-run form on Phoenix, quoted above, with
\texttt{eval/returns} flat for tens of millions of frames (\S\ref{sec:meas-collapse}). Second,
the critic at a collapse is not
uninformative, only quiet: an ordering with small gaps still ranks the actions, and uniform is
the one map that discards all of it. There is no intrinsic-motivation term in this loop
(\S\ref{sec:contracting}), so the critic is the only thing directing the actor, and a rule that
turns a small signal into no signal leaves nothing driving the run.

(A) fails: this is the Enduro failure and the observed collapse, and it is self-reinforcing.
(C) fails: $\rho^\star$ becomes infeasible as $\gapbar$ contracts under it, and the most diffuse
policy meeting an infeasible budget is the uniform one.

\textbf{(A$'$) and (C$'$) fail, and the normaliser is what makes them interesting anyway.}
$\snorm$ is estimated from returns in the buffer, not from the network, so it does not
contract when the network's outputs do: the gaps shrink under a fixed $\tau = c\,\snorm$ and
the policy goes uniform, exactly as (A) does. But that is also the only way to \emph{see} the
collapse --- \S\ref{sec:norm}'s third use --- because every rule that passes this case passes it
by being blind to the difference between a fine-grained game and a dead value function.

\textbf{(D) passes, and this is the one case it is built for and wins.} It is exactly
scale-invariant against the perturbation $\tilde{q} \to \varepsilon\tilde{q}$: every gap and
hence the target scales by $\varepsilon$, the solution scale moves to $\varepsilon\tau$, and
$\sm(\varepsilon\tilde{q}/\varepsilon\tau)$ is the \emph{same} policy. The loop's gain is zero.
It does not push back against a contracting value function; it declines to amplify the
contraction, which is enough that the collapse cannot compound.

\textbf{(F) passes, and it is the change of denominator that makes it do so.} Both $\tilde{q}$
and $\srank$ are outputs of the same network, so under $\tilde{q} \to \varepsilon\tilde{q}$ the
scale moves to $\varepsilon\tau$ and $\sm(\varepsilon\tilde{q}/\varepsilon\tau)$ is the
\emph{same} policy: (F) has (D)'s zero loop gain, exactly, and by the same argument. The previous
revision's $\dbar$ did not --- the parameter noise is not a function of the collapsing outputs,
so as the gaps contracted under it the ratio $\gapbar/\dbar$ went to zero and the branch asked
for a near-uniform policy, which is operationally fatal whatever its epistemic merit. Stating the
rule on a spread of the critic's own outputs is what collapses the two failure modes into one
denominator.

(I) passes on the same condition it carried at \S\ref{sec:c1} and for the same reason: an
epistemic spread of the network's own action values contracts with them and the loop gain is
zero, while a novelty score does not contract with them and its carrier decides the case. (E) is
exposed here in principle and much less in practice --- \S\ref{sec:meas-return} scores the
argmax arm well above every soft arm over exactly the frames where collapse lives, with the
identical parameter noise --- so
whatever the draw does early, adding an undirected softmax on top of it is what costs the budget.

\emph{The floor survives this reasoning, with a smaller job.} $\gapbar/k_0$ is no longer the only
scale-free thing in the rule and no longer carries \S\ref{sec:c6} on its own; it is a guard
against $\srank$ itself failing (\S\ref{sec:rule}), and \S\ref{sec:meas-k} sets it where it does
not bind.

\subsection{The horizon: regret rather than off-greedy rate}
\label{sec:c7}

\emph{Setup, at two ends.} Fix an off-greedy probability $p$ per step and ask what it costs over
an episode of $H$ steps.
\emph{(i)} On states where the actions are near-tied, so that each off-greedy action is taken at
a gap of essentially zero.
\emph{(ii)} On a state where one action must be repeated: action $R$ pays $1$ per step, action
$N$ pays nothing, and the episode lasts $T$ steps, so the optimal return is $T$ and a policy that
takes $R$ with probability $1-p$ scores $(1-p)T$. With no time limit both policies diverge and
the two separate only under discounting, where the gap is $\tilde{q}(s,R) - \tilde{q}(s,N) = 1$
exactly.

\emph{Demand.} The quantity held fixed across states may not be the off-greedy rate. In (i) the
rate is nearly free, so a fixed one is a tax that buys nothing wherever the state is already
resolved; in (ii) the same rate is paid in full and linearly in the horizon, and the agent scores
a fixed fraction of the maximum forever. No single $p$ is right for both. What separates them is
\emph{not} the horizon: the cost at either end is $p\,\gap\,H$, and $H$ is common to every state,
so it cancels from any comparison between them. What does not cancel is the price of one
off-greedy action, which is the gap. So the currency is regret, $p$ weighted by what it costs,
and a rule that controls $p$ directly is controlling the wrong number.

\emph{(i) retires a criterion, and it is the criterion the horizon tempts one into.} If what
mattered were executing an intended trajectory intact, the quantity to hold would be $(1-p)^{H}$,
and at any $H$ of this size that is negligible for every $p$ an exploring agent could want: an
agent that requires clean trajectories cannot explore at all. \S\ref{sec:meas-rho} confirms the
premise --- no arm this project has run comes close to a clean trajectory, and all of them work
--- so trajectory purity is not the currency. Regret is, and the same measurement prices it:
$\epsilon$-greedy costs an order of magnitude more than a softmax at a comparable rate, because
$\epsilon$ selects uniformly among exactly the actions the value function holds unlikely.

\emph{(ii) is Phoenix's own shape}, a game whose optimal policy is to keep firing, and it is
carried for one reason: it bounds (i)'s reading. ``The rate is nearly free'' is a measured fact
about the states this agent visits and not a general one, and a rule calibrated to a rate that is
free at a near-tie will pay that same rate at a state where it is worth $p\,T$.

\emph{What (ii) separates, and it is a separation of form rather than of degree.} A Boltzmann
rule puts $p_N = (1 + e^{\gap/\tau})^{-1}$ on the wrong action, which is \emph{exponentially
small} in $\gap/\tau$; $\epsilon$-greedy puts $\epsilon(1-1/|\actions|)$ there, which does not
depend on $\gap$ at all. At a state whose gap is large in units of the temperature the two are
not close, and they are not close by a factor that grows without bound as the gap does. That is
the whole of it: \textbf{a rate that answers to the gap against a rate that does not}, and (ii)
is the construction that makes the difference unbounded rather than merely large.

\eqref{eq:rule} collapses here for a second and independent reason worth stating, because it is
exact rather than exponential. The returns in (ii) are deterministic, so every action's
distribution is a point mass, hence a translate of one shape; by \S\ref{sec:bg}'s zero set
$\srank = 0$, so $\tau = 0$ and the rule is argmax at this state \emph{by construction}, paying
nothing at all. (E) collapses too, the gap being large against the parameter draw's own spread.
(C) does pay, at $p_N = \rho^\star/\gap$, but that is the budget it was handed rather than a
failure of this case.

So (ii) separates exactly one candidate, and it is the cleanest of the three places this document
retires it: no large $M$ as in \S\ref{sec:c2}, no measurement as in (i), only a gap and a horizon
to pay it over. What (ii) does \emph{not} do is separate any two \emph{survivors} --- all of them
are Boltzmann or argmax and all of them lose exponentially little here --- so it is no help in
choosing a denominator, which is this section's other job. Whether the loss falls once the
ordering has been resolved is \S\ref{sec:c11}'s question and not this one.

\textbf{This retires (B).} $\epsilon$-greedy prices no action, so it pays a rate the gap cannot
talk it out of: $\epsilon(1-1/|\actions|)T$ of the score in (ii), and in (i) an order of magnitude
more regret than a softmax at a matched rate (\S\ref{sec:meas-rho}). \S\ref{sec:c2} is the same
fact in its extreme form. It is also why
every surviving rule is a Boltzmann policy: value-weighting is the cheap way to buy a rate.
Everything else on the list passes, being either a Boltzmann policy or argmax.

\subsection{Exact ties}
\label{sec:c9}

\emph{Setup, in two variants, because they are not the same state.}
\emph{(i) Exact ties}: several actions have identical return distributions --- a no-op
duplicated, a direction the game ignores in this state.
\emph{(ii) Mean ties}: the actions have equal means and different distributions --- one is a
gamble and another is certain --- and the critic has resolved that this is so.

\emph{Demand.} The rule is well defined in both, and without a branch. At (i) every policy over
the tied actions is the same policy, so anything is correct and the only failure available is to
be undefined. At (ii) the means are equal but the actions are not, and the case asks the weaker
question of whether the rule notices.

(A) gives a uniform policy over tied actions in both, which is harmless. \textbf{(D) is
degenerate in both}: $\rho(\pi;s) = 0$ for \emph{every} policy, since $\rho$ reads the means, so
$\rho(\tau;s) = c\,\gapbar(s)$ reads $0 = 0$ and every $\tau$ in the bracket satisfies it. The
previous revision's bisection returns an arbitrary point, and under coupling that arbitrary
number multiplied the Munchausen log-policy. (C) and (C$'$) are infeasible rather than
degenerate, which is a softer failure with the same outcome.

(F) is well defined without a branch. At (i) identical actions have identical return
distributions, so $\srank = 0$ and $\tau = 0$ is argmax over tied values --- arbitrary, but
arbitrary among identical things, and no root-find exists to be degenerate. At (ii) $\srank > 0$
and the rule samples over them. (I) is also well defined and does the opposite at (ii): the case
says the critic has resolved the distributions, so there is no epistemic spread left and (I)
plays argmax. The case does not say which is right --- sampling among mean-tied actions collects
information about their distributions, and if that information is already in hand it buys nothing
--- and it is worth recording as the one place where (F) and (I) disagree about behaviour rather
than about price.

This is a small case with one large consequence: it is the second of two reasons (with
\S\ref{sec:rule}'s cost) to state the rule as a scale computed directly rather than as a
constraint solved numerically.

\subsection{Learned states with irreducible noise}
\label{sec:c10}

\emph{Setup.} The ordering at $s$ is settled and the network knows it, but the return from $s$
is genuinely variable --- sticky actions, a bonus that sometimes lands, anything the agent does
not control.

\emph{Demand.} The policy is sharp. Variance the agent cannot reduce is not a reason to try
other actions. This is the distinction between not knowing which action is best and knowing it
while the outcome stays uncertain, and every rule above that reads a spread is at risk of
confusing the two.

(A), (A$'$), (B), (C), (C$'$) and (D) pass by not looking: none of them reads a spread of the
return at all. (E) is exposed in principle --- $\snoise$ is fitted by gradient on a loss whose
residual is inflated by exactly this noise --- and there is no measurement either way. (I) passes
by construction, since an epistemic signal is defined as the part that contracts with data and
irreducible noise is the part that does not; that is the whole reason the family is carried.

\textbf{(F) passes for action-independent risk and fails for heteroscedastic risk}, and the
split is what chooses between the two free denominators. Risk the action does not control enters
$Z(s,\cdot)$ as a common term, which the dueling head puts in the value stream and the centred
advantage removes, so $\srank$ does not see it and $\sret$ sees nothing else. A gamble that one
particular action takes does move $\srank$, and there the rule diffuses when it should not. So
$\sret$ is retired as an operating denominator here and kept only as a unit (\S\ref{sec:norm}),
and $\srank$'s residual exposure is the honest price of a statistic that costs two reductions.

\emph{The exposure is wider than a gamble, and \S\ref{sec:bg}'s zero set says how much.} What
the centring removes is a common \emph{translation}, so what survives is every departure from a
location family --- not only an action whose outcome is a gamble, but any action whose return
has a different \emph{shape} from its neighbours', including one that dominates them at every
risk level. A settled ordering is therefore not sufficient for $\srank$ to be small. That the
rule nonetheless transfers (\S\ref{sec:denoms}) is a measurement and not a consequence of the
definition, which is the right way round for the claim most worth breaking.

Nothing free passes this case \emph{because} it separates the two uncertainties; (F) passes the
part of it that the architecture happens to separate for free. Doing it on purpose turns out to
cost two draws rather than a head --- \eqref{eq:uasplit} estimates both halves without bias from
two posterior samples of the same network --- so the price of this case is
Remark~\ref{rem:peraction}'s, not \S\ref{sec:contracting}'s.

\subsection{Repeated visits to a resolved state}
\label{sec:c11}

\emph{Setup.} A state visited over and over, its ordering unchanged, its returns collected.

\emph{Demand.} Sharpness rises. An exploration rate that does not anneal is a tax with no term,
and at a full budget the late fraction of the run is where most of the score is.

(A) and (A$'$) fail: a constant is a constant, and a constant in normalised units is a constant.
(B) anneals, but on a clock rather than on the state, and BTR's geometric decay is the reason its
arms are competitive early --- a schedule is a real answer to this case and it is the only thing
in the project that has ever given one. (C), (C$'$) and (D) fail: they read gaps, and the gaps at
a converged state stop moving.

\textbf{(E) and (F) fail, and this is measured rather than argued.} Across Phoenix checkpoints
spanning most of a budget, \S\ref{sec:denoms} finds both $\dbar$ and $\gapbar/\srank$ flat.
Neither
the parameter draw nor the quantile spread contracts with data on this codebase. The rule
therefore holds a constant \emph{relative} sharpness for a whole run, which is better than a
constant absolute one and is not annealing.

Only (I) passes, and only because contraction is what a posterior, a count or an ensemble's
disagreement is \emph{for}: they go as $1/\sqrt{n}$ where nothing above goes at all. Fox's
closed-form temperature (\S\ref{sec:related}) passes it by construction for the same reason, its
denominator being the epistemic variance of the gap, which falls as $1/n$.

What the case costs is decided by Remark~\ref{rem:peraction} rather than by this document. Given
draws that are actually posterior samples, the epistemic half $\tilde\sigma_{\mathrm{epi}}$ of
\eqref{eq:uasplit}'s two-draw split is a contracting per-state signal for one extra pass through the output layers. Given the draws as
they are wired, it is not, and then the case needs a head. \S\ref{sec:readmit} is ordered
accordingly.

\subsection{Summary}
\label{sec:summary}

\emph{(A)} fixed sharpness, the status quo; \emph{(A$'$)} the same in normalised units;
\emph{(B)} $\epsilon$-greedy; \emph{(C)} a fixed absolute regret budget; \emph{(C$'$)} the same,
normalised; \emph{(D)} a fixed fractional one, two revisions ago's recommendation; \emph{(E)}
argmax with the parameter draw, the best-measured arm; \emph{(F)} \eqref{eq:rule}; \emph{(I)} a
rule scaled by something that contracts with visitation.

\begin{center}
\small
\setlength{\tabcolsep}{4.5pt}
\begin{tabular}{lccccccccc}
  \toprule
  case & (A) & (A$'$) & (B) & (C) & (C$'$) & (D) & (E) & (F) & (I) \\
  \midrule
  \S\ref{sec:c1} rescaling        & \no  & \yes & \yes & \no  & \yes & \yes & \meh$^\dagger$ & \yes & \meh$^\ddagger$ \\
  \S\ref{sec:c3} duplicates       & \yes & \yes & \yes & \yes & \yes & \yes & \yes & \yes & \yes \\
  \S\ref{sec:c4} added actions    & \meh & \meh & \yes & \yes & \yes & \no  & \meh & \meh & \meh \\
  \S\ref{sec:c2} catastrophic     & \yes & \yes & \no  & \yes & \yes & \no  & \yes & \yes & \yes \\
  \S\ref{sec:c5} two rooms        & \no  & \no  & \meh & \meh & \meh & \meh & \meh & \yes & \yes \\
  \S\ref{sec:c6} uninformative    & \no  & \no  & \yes & \no  & \no  & \yes & \meh & \yes & \meh$^\ddagger$ \\
  \S\ref{sec:c7} horizon          & \yes & \yes & \no  & \yes & \yes & \yes & \yes & \yes & \yes \\
  \S\ref{sec:c9} ties             & \yes & \yes & \yes & \meh & \meh & \no  & \yes & \yes & \yes \\
  \S\ref{sec:c10} learned, noisy  & \yes & \yes & \yes & \yes & \yes & \yes & \meh & \meh & \yes \\
  \S\ref{sec:c11} revisitation    & \no  & \no  & \meh & \no  & \no  & \no  & \no  & \no  & \yes \\
  \midrule
  constants with a value scale    & $1$ & $0$ & $0$ & $1$ & $0$ & $0$ & $0$ & $0$ & $0^\ddagger$ \\
  extra parameters                & $0$ & $0$ & $0$ & $0$ & $0$ & $0$ & $0$ & $0$ & ${>}0$ \\
  extra work per step             & --- & scale & --- & solve & both & solve & --- & 2 red. & a head \\
  \bottomrule
\end{tabular}
\end{center}

$^\dagger$ conditional on $\snoise$ tracking the value scale --- measured not to,
\S\ref{sec:denoms}, though (E) fails toward sharpness. $^\ddagger$ a contracting signal is
generally not in value units, so (I) needs a carrier and is scale-free and collapse-neutral only
with one; \S\ref{sec:contracting}. ``scale'' is a running return normaliser, ``solve'' a
per-state root-find, ``2 red.'' two reductions over an array already formed. (G), an entropy
target, is excluded at \S\ref{sec:c3} and is not carried.

Read down the columns. (A) is the status quo and fails three cases including both that this
project has observed. (A$'$) buys exactly one of them back and is the measure of what a global
normaliser is worth on its own: it is the difference between a constant that is wrong between
games and a constant that is wrong between states. (B) fails the two that are about pricing.
(C) and (D) fail complementary halves and neither dominates the other --- which is why the
previous revision's move from (C) to (D) fixed \S\ref{sec:c6} and broke \S\ref{sec:c2}; (C$'$)
repairs (C)'s rescaling failure and leaves the rest of its column standing. (E) is the
best-measured arm and fails nothing outright except the last case, which is a result in itself
and is why \eqref{eq:rule} is built to contain it. (F) is what is left of what is free. The last
two rows of cases are what is not free, and the last column is the only thing in the table that
answers them.

\section{Coupling the temperature into the loss}
\label{sec:decouple}

The question is whether $\tau(s)$ replaces $\tau_\ell$ in the Munchausen target as well as in
\texttt{\_act}. The previous revision answered yes on principle --- M-DQN's output \emph{is} a
policy, $\sm(\tilde{q}/\tau_\ell)$ exactly, so acting at any other scale deploys something the
algorithm did not produce --- and the revision after it answered no, on an objection to the
objective a state-dependent coefficient induces. That objection is now measured, and it does not
hold (\S\ref{sec:c8}).

What follows separates two questions the previous revisions ran together. \emph{Is coupling the
better-posed design?} Yes, and \S\ref{sec:couple-case} says so in four parts, because the rest of
this section is a list of obstacles and none of them is a defence of the fixed constant.
\emph{Is it the thing to run next?} No --- for one reason that is about what an experiment could
conclude and not about the design (\S\ref{sec:c8-staging}). Between the two sits the one obstacle
that is neither, and it is the only place in this section where something is actually unknown:
the fixed point (\S\ref{sec:couple-fixedpoint}).

\subsection{The case for coupling}
\label{sec:couple-case}

\paragraph{The critic's output is a policy at exactly one scale, and decoupling is not playing
it.} By Proposition~\ref{prop:amp}, $\sm(q_\infty/\tau_\ell)$ is the soft-optimal policy of the
MDP the loss is solving. Decoupled, the behaviour is $\sm(\tilde{q}/\tau(s))$, which is that
policy at no scale the loss knows about: the agent solves one problem and plays another, and the
mismatch is a function of the state. Coupled, the identity is restored state by state and by
construction --- $\tau(s)$ is the loss's own coefficient at $s$, so the behaviour is the
soft-optimal policy of the MDP the critic is being fitted to. That is the previous revision's
principle in full, and nothing measured since has touched it.

\paragraph{It deletes a constant rather than adding one.} Decoupled, the rule needs the ceiling
$\tau(s) \le \tau_\ell$, and Remark~\ref{rem:erratum} is what makes that bound the right one
rather than an arbitrary one: $\tau_\ell$ is the scale at which the critic is a policy. Coupled,
there is no $\tau_\ell$ left to be a ceiling --- the loss's coefficient \emph{is} $\tau(s)$, the
critic is a policy at $\tau(s)$ by the same proposition, and the bound has nothing left to say.
So coupling removes \texttt{MUNCHAUSEN\_TEMPERATURE} from the algorithm outright, and with it the
one constant in the loss that carries a value scale --- which is \S\ref{sec:project}'s standing
requirement, applied to the one place this document had exempted from it. \emph{And the bound it
removes is a structurally fragile one.} The fraction of states resting on the ceiling is the
fraction with $\srank(s)/g > \tau_\ell$, so it is a quantile of the $\srank$ distribution read at
a threshold set by the loss --- and a quantile taken in the body of a distribution moves far
faster than the threshold does. \textbf{So the decoupled rule's resting fraction is hostage to a
constant chosen for the loss's reasons, and not gently}: \S\ref{sec:meas-k} measures a modest
change in $\tau_\ell$ turning the ceiling from a guard into the branch that decides a large
minority of actions --- over which \eqref{eq:rule} is the fixed temperature it exists to replace.
Coupled, the branch does not exist.

\paragraph{The distortion it induces is smaller than the one decoupling already accepts.}
Both are per-step additive terms in $\tilde{q}$'s units --- $D(s)$ on one side,
\S\ref{sec:decouple-ledger}'s lost Munchausen cancellation on the other --- so they compare
directly, and \S\ref{sec:meas-couple} finds $D$ smaller by two orders. In the loss's own units
there is no argument for the decoupled arm. \emph{That comparison is between two terms in the
loss and is not a claim that the bonus is behaviourally negligible}; read instead as a standing
bias on the policy, at the per-decision horizon, the same $D$ is of the same order as this
project's existing exploration mechanism, which is \S\ref{sec:c8}'s second bullet and the reading
that bears on where a coupled agent would go.

\paragraph{And it makes the Munchausen clip inert by identity rather than by luck.} The bonus is
$\alpha\tau\ln\sm(\tilde{q}/\tau)(a) = -\alpha(\gap_a(s) + D(s))$ identically, by
\eqref{eq:distortion}'s algebra, so $l_0 = -1$ fires iff $\gap_a + D > 1/\alpha = 1.111$ --- a
condition on gaps in value units, with the temperature cancelled out of it --- so counted over
state-action pairs the clip is \emph{not} the dead letter the previous revision recorded, and
\S\ref{sec:meas-couple} says how far from dead. But the loss evaluates the bonus only at the
action the behaviour took, and there the two conditions collapse into one: the clip needs
$\pi_\ell(a) < e^{-1/(\alpha\tau_\ell)}$, and when the behaviour shares that temperature this is
\emph{also} the probability of drawing $a$. \textbf{Coupled, the chance a drawn transition is
clipped is therefore at most $|\actions|\,e^{-1/(\alpha\tau(s))}$ at every state}, a bound that
needs nothing measured and is minute at any sharp $\tau$. Decoupled there is no such identity:
the rate runs through the behaviour's own temperature,
$\pi_b(a) = \pi_\ell(a)^{\tau_\ell/\tau_b}$, which is safe while $\tau(s) \le \tau_\ell$ and
degrades as a power of the mismatch above it. The clip is an argument for coupling, and the
previous two revisions had it backwards in both directions at once.

\emph{So there is no argument against coupling from the objective, from the constants, or from
the clip.} The three that remain are below, and only the first is about the algorithm.

\subsection{The objective under a state-dependent coefficient}
\label{sec:c8}

M-VI($\alpha, \tau$) solves the MDP regularised by entropy at $\lambda = (1-\alpha)\tau$, so a
state-dependent $\tau(s)$ solves one regularised at $\lambda(s) = (1-\alpha)\tau(s)$, and the
agent's objective gains \eqref{eq:distortion} at every step: \textbf{states with a large $D(s)$
are literally worth more}. That is not a change of temperature, it is an intrinsic reward, and
the previous two revisions both stopped there --- they named the reward's \emph{character} from
the denominator and ranked the candidates on that. Character without magnitude decides nothing,
and both are now measured (\S\ref{sec:meas-couple}). Three readings, and each moves a claim.

\paragraph{What the bonus is paid on is the near-tie, not the denominator.} $D$ is linear in
$\tau$ and exponential in $-\gap/\tau$, so the exponential chooses the states and the denominator
only scales what they get. Weighting each per-state quantity by the bonus it carries,
\S\ref{sec:meas-couple} finds the mass sitting on states whose top-two gap is far below average,
while $\srank$'s own concentration is slight on one game and real on the other. So the intrinsic
reward pays the agent to stand where \emph{the two best actions are nearly tied in mean}, which is
very nearly the opposite of what the objection charged it with, and which is the cheapest place in
the state space to be paid to stand: the regret of being wrong there is small by construction. The
$\srank$ carrier the objection was built on is the weaker of the two wherever both are present.

\emph{This retires the risk-seeking charge and does not retire the case behind it.}
\S\ref{sec:c10} stands as a case about the \emph{sampling} rule --- a heteroscedastic state gets a
diffuse policy, which is a real failure of (F) and is why \S\ref{sec:c10} is in the table. What
does not follow is that coupling converts it into a risk-seeking intrinsic reward, because the
conversion runs through an exponential that ignores $\srank$ almost entirely.

\paragraph{The magnitude is orders below the cost of not coupling.} Both $D$ and the one lost
Munchausen cancellation \S\ref{sec:decouple-ledger} accepts are per-step additive terms in
$\tilde{q}$'s units, so they compare directly; \S\ref{sec:meas-couple} makes that comparison on a
game where both were measured, and it goes the wrong way for the objection this subsection exists
to state.

\emph{And the ranking against the previous revision's rule was inverted.} $D$ grows with how
diffuse a rule is, and \eqref{eq:rule} is calibrated to be sharp where candidate (D) is
calibrated to be diffuse, so priced on the same states (D)'s induced bonus is the larger by two
orders --- large enough to match the cancellation \S\ref{sec:decouple-ledger} treats as a real
cost (\S\ref{sec:meas-couple}). The bullet that condemned this denominator names the one rule in
the family whose induced bonus is negligible, \emph{because} it is the sharpest, while the rule
it ranked above sits far above it. The objection was sound and it was attached to the wrong
candidate.

\paragraph{The channel is gated by the gap, which is the finding that matters.} Whatever a
coupled rule divides by, it reaches the objective only through \eqref{eq:distortion}, and
$\h(\pi)$ multiplies it. Push $\tau$ all the way to $\tau_\ell$ at every state --- the most a
coupled rule bounded that way can ever pay --- and \S\ref{sec:meas-couple} finds the great
majority of states still receiving under a tenth of the ceiling $\tau_\ell\ln|\actions|$, because
their gaps are large in units of $\tau_\ell$; under \eqref{eq:rule} itself, nearly all of them.
\textbf{A state that is genuinely unresolved but
where the critic is confidently wrong receives nothing.} A reward-level bonus --- RND's, a
count's --- is not gated that way; it is added whatever the critic believes. This is what
separates the entropy channel from an intrinsic reward proper, and it is the objection
\S\ref{sec:coupling} has to answer.

\emph{So the verdict on the three denominators, restated on what was measured.}
\begin{itemize}
  \item \emph{With $\tau \propto \gapbar$} --- the previous revision's rule --- $k$ is constant
  by construction, so $\h$ is roughly state-independent and the bonus really is proportional to
  the denominator, which \S\ref{sec:meas-couple} confirms: the mass tracks $\gapbar$ and is
  spread rather than concentrated. The agent is paid to occupy states with a large value spread,
  and by \S\ref{sec:c2} those are the ones containing the worst actions. This is the one
  denominator for which the character argument goes through, it is confirmed rather than argued,
  and its magnitude is too large to rescue it.
  \item \emph{With $\tau \propto \srank$} --- this document's rule --- substitute the branch into
  \eqref{eq:distortion} and the answer is closed-form rather than a matter of character. On
  $\tau = \srank/g$, which \S\ref{sec:meas-couple} finds carries most of $\sum D$,
  \[
    D(s) \;=\; \frac{\srank(s)}{g}\,\ln m(s),
    \qquad
    m(s) \;=\; \sum_a \exp\!\left(-\,g\,\frac{\gap_a(s)}{\srank(s)}\right),
  \]
  exactly. The bonus is \emph{linear} in $\srank$ and, through $m$, \emph{exponential} in every
  gap measured in units of $\srank$ --- so what it responds to is each action's gap divided by
  the state's own ordering spread, and not either quantity alone. Where a single runner-up is in
  contention this collapses to $(\srank/g)\exp(-g\,\gap_2/\srank)$ and \textbf{one ratio decides
  it}, $\gap_2/\srank$, on which \S\ref{sec:meas-couple} confirms the bonus concentrates far more
  sharply than on either part of it. Where several are in contention it does not collapse, and
  \S\ref{sec:meas-couple} measures how often that matters.

  \textbf{So the agent is paid to occupy states where the top two actions are close in mean
  \emph{relative to how far their ordering swings across the return distribution}} --- where the
  critic's ranking of the best two is not robust to which part of the distribution it is read at.
  That is a sharper object than a near-tie, and it excludes the obvious degenerate case: a perfect
  tie between actions whose returns are translates of one shape has $\srank = 0$, hence $\tau = 0$
  and $D = 0$, and receives nothing. What is rewarded is a tie the return distribution itself
  disputes.

  \emph{Two consequences follow, and they are what \S\ref{sec:coupling} has to weigh.} The bonus
  is delivered through the backup, so it is discounted and propagated: the agent does not merely
  stand in such states, it comes to value the states that \emph{lead} to them. And $\srank$ is
  measured aleatoric (\S\ref{sec:c10}) and does not contract with data (\S\ref{sec:c11}), so
  \textbf{the bonus never extinguishes}. Visiting one of these states resolves nothing about it,
  because there is nothing there to resolve --- which action is better genuinely depends on the
  risk level. This is the difference between it and family (I) below, and it is not a difference
  of magnitude: it is a permanent preference for risk-ambiguous states rather than an exploration
  bonus that spends itself.

  \emph{Two currencies, and they do not agree.} As a per-step term in the loss $D$ is minute
  against what decoupling already costs. But $D$ is a per-step \emph{value} term, which is exactly
  what \S\ref{sec:meas-rho} prices the exploration mechanisms in, and over a horizon $H\,D/Q^\star$
  is what a standing bias on the policy is worth --- a quantity larger by the factor $H$ and, as
  \S\ref{sec:meas-couple} measures it, of the same order as this project's entire existing
  exploration mechanism. Both readings are correct and they answer different questions: the first
  compares two terms in the loss, the second asks how far the bonus moves the policy. Only the
  second answers where the agent would go, and the previous revisions quoted only the first.

  \emph{And the risk-seeking charge is reduced rather than retired.} $\sret$ --- the width of the
  return, which is what the charge was actually about and which no previous revision measured ---
  does carry some concentration, slight on one game and real on the other
  (\S\ref{sec:meas-couple}). So the charge was aimed at the wrong quantity on one game and at a
  weakened version of the right one on the other.

  \emph{And an aliased action set revives it outright, which is the sharpest thing against
  coupling this denominator.} By \S\ref{sec:bg}, duplicating actions sends
  $D \mapsto D + \tau\ln(\text{copies})$, and on this branch $\tau = \srank/g$, so the added term
  is \textbf{exactly proportional to $\srank$} --- the pure $\srank$ bonus the objection
  described, arriving not from the near-tie mechanism but from how many ways the environment
  happens to spell the same action. \S\ref{sec:c3} is the case that the \emph{rule} survives this
  and it does; the coupled \emph{objective} does not. Two things follow. An action set with
  aliases pays a standing bonus proportional to the denominator wherever the aliased actions are
  in contention, and \S\ref{sec:c3} already identifies one of this project's two games as heavily
  aliased --- which is also, and perhaps not by coincidence, the game whose measured $\srank$
  concentration is the anomalous one (\S\ref{sec:meas-couple}). And the term is not a small
  correction: \S\ref{sec:meas-couple} prices a single duplication at several times the entire
  bonus it perturbs. \textbf{So the character of a coupled $\srank$ rule is not a property of the
  rule alone; it is a property of the rule and the action set together}, and no measurement on one
  action set settles it for another. That is a reason to want the contracting denominator of
  family (I), whose induced term carries the same duplication factor but spends it.
  \item \emph{With $\tau \propto \hat{s}$, a contracting signal} --- the family (I) --- the
  induced term is propagated exactly as \eqref{eq:rule}'s is, and differs in the one way that
  matters: it \emph{falls} as the ordering resolves, so the bonus is spent rather than drawn
  forever. That is the property separating a bonus from dithering, and it is what the bullet
  above lacks --- not the propagation, which coupling gives either denominator, but the
  extinction. But it is delivered through the same gated channel,
  so it is an exploration bonus that arrives only where the critic is already indifferent. That
  is a weaker instrument than \S\ref{sec:coupling} assumed, it is measurable before any run is
  spent, and it is now \S\ref{sec:coupling}'s first item rather than its premise.
\end{itemize}

\begin{remark}[The contraction result, and what it does not cover]
\label{rem:coupled-ok}
A state-dependent coefficient breaks nothing so long as it is fixed in advance:
$T_\Omega q = r + \gamma P \Lambda_\Omega q$ with $(\Lambda_\Omega q)(s) = \lambda(s)\lse(q(s,\cdot)
/\lambda(s))$ is a $\gamma$-contraction for any $\lambda(s) > 0$, because the argument is pointwise
in $s$ and regularised-MDP theory asks $\Omega_s$ to be strongly convex at each state and not to
be the same at every state; and M-VI's fixed point is
$q_\star = A q_*^{\lambda} - \alpha A v_*^{\lambda}$ uniquely --- $q_*^{\lambda}$ and
$v_*^{\lambda}$ being the action- and state-values of the MDP regularised at $\lambda$ ---
with $A = 1/(1-\alpha)$, gap amplification $A$ per state and
$\sm(q_\star(s,\cdot)/\tau(s)) = \pi_*^{\lambda(s)}$.
\texttt{scripts/check\_mvi\_gap.py} confirms both to six decimals with $\tau(s)$ spanning a decade.
\textbf{Both results take $\lambda$ as exogenous}, and a rule that reads $\lambda$ off the
iterate is not covered by either; \S\ref{sec:couple-fixedpoint} is what that costs and what buys
it back.
\end{remark}

\subsection{The fixed point, once the coefficient depends on the iterate}
\label{sec:couple-fixedpoint}

This is the only obstacle in this section that is about the algorithm rather than about the
experiment, and it is the reason a coupled arm cannot simply be started.

\emph{What is not covered.} Remark~\ref{rem:coupled-ok}'s script draws its per-state temperatures
once from a fixed vector and holds them. A coupled rule sets
$\lambda(s) = (1-\alpha)\srank(q(s,\cdot))/g$, a functional of the very iterate being updated.
That operator is nonlinear and self-referential and is not the one analysed: pointwise strong
convexity is a statement about a fixed $\Omega_s$ and says nothing when $\Omega_s$ moves with the
iterate. There is also a feedback loop whose sign nobody has checked --- a larger $\lambda(s)$
smooths the values at $s$, which moves $\srank(s)$, which moves $\lambda(s)$. It is confined to
one state, so it does not compound through the discount, but that is an observation and not an
argument.

\emph{The cheap sufficient condition, and by how much it misses.} $\lambda\lse(q/\lambda)$ is
nonexpansive in $q$ in sup norm, and its derivative in $\lambda$ is $\h(\pi) \le \ln|\actions|$,
so for two iterates
\begin{equation}
  \|Tq_1 - Tq_2\|_\infty \;\le\; \gamma\left(1 + \frac{L(1-\alpha)\ln|\actions|}{g}\right)
  \|q_1 - q_2\|_\infty,
  \label{eq:couplebound}
\end{equation}
with $L$ the Lipschitz constant of $\srank$ in the same norm --- at most $2$ by inspection, a
centring and a spread each nonexpansive, and never measured. A contraction therefore wants
\begin{equation}
  g \;>\; \frac{\gamma\,L\,(1-\alpha)\ln|\actions|}{1-\gamma} \;=\; \gamma L(1-\alpha)H\ln|\actions|,
  \label{eq:couplegain}
\end{equation}
and the shape of that condition is the whole point: \textbf{it is proportional to the horizon}.
The same $H$ that \S\ref{sec:c7} treats as the reason to price exploration in regret is what
amplifies a per-state perturbation of the backup here, so at a discount this close to one the
sufficient gain is orders above anything a behaviour policy would choose --- \S\ref{sec:meas-g}
calibrates $g$ and \S\ref{sec:meas-couple} evaluates \eqref{eq:couplegain} against it, and they
do not meet. The bound is crude: it takes the worst case over states and over $\lambda$ at once,
and it is linear in $1/g$ where the actual sensitivity is exponentially gated (\S\ref{sec:c8}).
Failing it is not evidence that the scheme diverges. It is evidence that the cheap argument is
not available, so something has to be run.
A further gap is worth naming: $\srank$ is a functional of the return \emph{distribution} and
not of $q$, so even \eqref{eq:couplebound} is an analogy --- the distributional operator
contracts in a Wasserstein metric, not in sup norm on the means.

\emph{Two things clear it, either sufficient, both cheap.}

\textbf{Read $\tau$ off the target network.} Then $\lambda$ is a fixed function of $s$ between
synchronisations, every period is exactly Remark~\ref{rem:coupled-ok}'s operator, and the coupled
scheme is no worse founded than the target network the method already rests on. $\tau(s')$ is
free --- \texttt{target\_next\_qs} is already in registers in \texttt{\_train\_step}. $\tau(s)$ is
not: the Munchausen log-policy is read off the \emph{online} pass, deliberately, and a target
pass over \texttt{obs} is a whole extra trunk evaluation, whose throughput cost
\texttt{btr.py}'s performance block already measures. That is the honest price of the clean
version, and it is the first thing to trade against staying decoupled.

\textbf{Iterate it tabularly first.} \texttt{scripts/check\_mvi\_gap.py} already builds random
finite MDPs and iterates M-VI against entropy-regularised VI; the change is to recompute
$\lambda(s)$ from the iterate on every sweep instead of drawing it once. Whether the composed map
converges, to what, and with what dependence on $g$ then stops being a theorem nobody has and
becomes a measurement that costs minutes. \textbf{If the loss is to carry $\tau(s)$, this is the
first thing to do} --- before the target-net question, and before any arm. Measuring $L$ is the
same size of job and makes \eqref{eq:couplebound} quantitative rather than indicative: perturb
the head's output over replay states and read the induced $\Delta\srank$.

\subsection{Containment of the sharp limit}
\label{sec:c8-staging}

\eqref{eq:rule} is built to contain arm (E): as $g \to \infty$, $\tau(s) \to 0$ and the rule is
argmax acting on the parameter draw, which is the best thing this project has run
(\S\ref{sec:meas-return}). That containment is what makes a negative result cost one run and no
baseline, and it is the whole staging (\S\ref{sec:staging}). It is a property of the
\emph{experiment}, not of the rule, and this subsection is the only thing in the section that
turns on the difference.

\textbf{Coupling destroys it.} Coupled, $g$ is a coefficient in the loss, so $g \to \infty$ sends
$\lambda(s) \to 0$ and deletes the entropy regularisation the algorithm is defined by: the limit
is no longer arm (E) but M-DQN at zero temperature, whose bonus has saturated at $-\alpha\gap_a$
everywhere and whose target is a hard max --- a different algorithm that has never been run. The
full-budget references stop being the comparison at every $g$, not only at the interesting ones,
and
$\tau \to 0$ --- which \S\ref{sec:meas-return} and \S\ref{sec:meas-g} both prefer --- becomes
inadmissible rather than the limit the rule is designed around. Decoupled, the control path stays
bit-for-bit and the change is a diff in \texttt{\_act} that reverts by deletion
(\S\ref{sec:project}).

\emph{One run has tested the premise, at a constant temperature.} Exactly one arm this project
has run moved \texttt{MUNCHAUSEN\_TEMPERATURE} at all, with the acting temperature aliased to it,
and \S\ref{sec:meas-return} reports what happened: it leaves the family the other soft arms
belong to, in the specific sense that the frame shift which overlays each of them onto the
argmax-only curve does not overlay it, and it finishes the weakest of them. One seed, two factors
moved at once, and a constant is not \eqref{eq:rule}, so it is not evidence about coupling a
state-dependent rule. It is evidence for this subsection's premise: \textbf{touching the loss's
temperature moved the run off the curve the references describe, on the first and only attempt.}

\subsection{Four further obstacles}
\label{sec:couple-obstacles}

\paragraph{Both constants were calibrated on decoupled nets.} $g \approx 4$ and $k_0 \approx 20$
are read off the gaps of four decoupled runs (\S\ref{sec:meas-g}, \S\ref{sec:meas-k}), and
\S\ref{sec:meas-couple}'s figures are a linearisation around a point a coupled scheme would
leave, since the gaps are exactly what $\tau(s)$ in the loss would move. Nothing offline prices
that. The branch fractions --- which say whether the rule is operating or resting on a bound ---
are the first thing to log, under either design.

\paragraph{$\tau(s)$ has to be stop-gradiented by hand.} In \texttt{\_train\_step} the Munchausen
bonus is already outside the gradient, but the \texttt{stop\_gradient} wraps
\texttt{action\_log\_probs} and not the temperature multiplying it. A coupled $\tau(s)$ read off
the online pass enters \emph{outside} that wrapper, and left there it hands the network a
gradient path into its own regularisation coefficient: the loss can then be reduced by moving
$\srank$ instead of by fitting the return. Reading $\tau$ off the target network
(\S\ref{sec:couple-fixedpoint}) closes this as a side effect; the online reading needs the
wrapper widened deliberately. It is one line, and nothing fails loudly if it is missed.

\paragraph{The loss would estimate $\srank$ from a quarter as many quantile samples.}
The loss draws fewer quantile samples per state than \texttt{\_act} does, so a coupled
$\hat\tau(s)$ carries the larger relative sampling error --- of order $1/\sqrt{2(n-1)}$ in the
sample count $n$, which at this project's two settings is about a factor of two
(\S\ref{sec:meas-couple}). What that costs is bounded by the same gating as
everything else here: $\lambda\lse(q/\lambda)$ is convex in $\lambda$ with derivative $\h(\pi)$,
so noise in $\hat\tau$ biases the backup upward by about $\frac12\Var(\hat\tau)\,d\h/d\lambda$,
and $\h$ is exponentially small wherever the gaps are large in units of $\tau$. It is second
order in the sharp regime and it is not second order at the near-ties --- which are exactly the
states $D$ is paid on. The existing probe measures it at both sample counts.

\paragraph{PER is untouched either way.} Coupling changes the value of a TD error, not how many
there are or how they are drawn: one loss, one priority per transition, one sampling
distribution. The previous revision listed this as something decoupling buys. It buys nothing
here.

\subsection{What decoupling costs}
\label{sec:decouple-ledger}

\emph{Costs.} The behaviour is not the soft-optimal policy of the MDP being solved. The
measurable residue is one term: the Munchausen bonus is evaluated at the taken action, so a
behaviour the loss did not anticipate loses one cancellation, clipped into $[-\alpha, 0]$ and not
propagated through the discounted sum. \S\ref{sec:meas-couple} measures it across a range of
behaviour entropies and finds it a fraction of a per cent of \texttt{train/q} --- small, but two
orders larger than the $D$ that coupling would induce in its place, which is the comparison
\S\ref{sec:c8} turns on.

Direction matters and bounds it further: acting \emph{sharper} than the loss expects is
conservative, because the value function anticipates randomness the behaviour will not have;
acting more diffusely is the Enduro direction. So enforce $\tau(s) \le \tau_\ell$, which is the
correct ceiling and not an arbitrary one because the critic's output over $\tau_\ell$ is already
the soft-optimal policy of the MDP being solved (Remark~\ref{rem:erratum}), and log the resting
fraction.

\emph{The ceiling is slack at the constant in force, and slack by less than it looks.} Where it
binds the direction is favourable --- $\tau(s) = \tau_\ell$ exactly, the behaviour \emph{is} the
loss's policy and this ledger's one term is zero there --- so the cost is not that the branch is
wrong but that it is not the rule, and how much of the time it is not the rule is set by a
constant this document does not control. \S\ref{sec:meas-k} measures both how slack it is and how
fast that slack disappears when $\tau_\ell$ moves, and finds a higher gain the cheapest fix.

\emph{Buys.} \S\ref{sec:c8-staging}, and that is the whole of it. The previous revision credited
decoupling with two further things and neither survives. $l_0$ stays inert because the acting
temperature is at or below the loss's, not because the two are kept apart --- coupling delivers
the same inertness as an identity (\S\ref{sec:couple-case}) --- and the clip rate the previous
revision projected for its own coupled rule was counted over actions uniformly, which is not the
distribution the loss draws from (\S\ref{sec:meas-couple}). And PER is untouched under either
design (\S\ref{sec:couple-obstacles}).

\paragraph{Where this leaves it.} On the design, coupling is better posed and this section finds
nothing against it: the distortion is orders below the term decoupling accepts, it
removes a constant rather than adding one, and it makes the clip's inertness an identity. What
stands between it and a run is \S\ref{sec:couple-fixedpoint} and \S\ref{sec:c8-staging} --- one
open question about the operator, which two cheap checks would settle, and one about what a first
experiment could conclude. \textbf{The order those imply is to decouple the first arm, because it
is the one that prices $g$ against references that still apply, and to settle
\S\ref{sec:couple-fixedpoint}'s tabular check while it runs.} If the composed map converges and
the target-net reading is taken, the coupled arm is the one to run second, with $g$ measured
rather than assumed. That is a claim about sequencing. Nothing in this section defends
$\tau_\ell$ as a constant, and \S\ref{sec:couple-case} is the reason the sequence ends with it
gone.

\section{Removals from the previous revision}
\label{sec:dropped}

The previous revision carried four proposals and a supporting apparatus. None of it is in
\eqref{eq:rule}; \S\ref{sec:readmit} takes some of it back on a different account, which is the
last two cases rather than the choice of temperature. In descending order of what it cost:

\emph{A fixed grid of $K$ members with a margin criterion}, at $2K$ heads and a parameter cost
parameters. It was the only proposal with a global guarantee --- within $c$ of the best member of
the grid --- and it needed a placement that nothing supplies, being a set of constants in value
units, i.e.\ \S\ref{sec:project}'s second constraint $K$ times over. Retired by \S\ref{sec:meas-return}'s four-arm table
regardless: it is an apparatus for choosing a level that is worth less than a seed.

\emph{A servo on the marginal value of an octave}, at four heads, a probe policy, a
difference-parameterisation to keep two output layers calibrated to within a tenth of a per cent,
and a concavity assumption that the fixed-weight table does not support. Same retirement, plus its
own: a servo sees a slope, so a non-monotone $J(\tau)$ --- the return of the policy at acting
scale $\tau$, \S\ref{sec:meas-return} --- gives it several fixed points and it tracks whichever
branch it started on.

\emph{A servo on exploration regret}, at one scalar head. The cheapest of the three and the only
one with a monotonicity condition that could be logged rather than assumed. Retired for the same
reason as the other two, and separately by \S\ref{sec:meas-g}'s arithmetic: without a second value
head to price the parameter draw, the head-free version saturates on a floor that is worst early
and decays with $\snoise$, making it a controller whose authority is scheduled by an unrelated
parameter.

\emph{The fixed fractional-regret rule} --- the previous recommendation --- retired by
\S\ref{sec:c2}, \S\ref{sec:c4} and \S\ref{sec:c9}. What survives of it is \eqref{eq:rule}'s floor
branch and the reason for it, \S\ref{sec:c6}, which is a real case that it is the only prior
candidate to pass.

\emph{The apparatus each of those needed} --- a $\tau$-conditioned network, the residual
parameterisation, per-member normalisers, the PER reference-member problem, the
leak-versus-deadband gate, the $n$-step correction for a per-member $Q_\env$, and the
bisection inversion of the regret --- was in every case a cost of holding several members,
of servoing a level, or of solving a constraint numerically, and \eqref{eq:rule} does none of the
three. \textbf{No normaliser is needed to make its constants portable}; that is narrower than the
previous revision's claim, and \S\ref{sec:norm} says what a normaliser is still for.

\emph{One thing is worth keeping and is not part of the rule.} A scalar $Q_\env$ head ---
\eqref{eq:qenv} below, one \texttt{stop\_gradient} at the trunk's activations, a negligible
parameters, off the control path and reversible --- buys three measurements available no other
way: the amplification ratio $\gapop(Q_\tot)/\gapop(Q_\env)$, which
Proposition~\ref{prop:amp} predicts settles near $1/(1-\alpha) = 10$ and which is that theory's
only observable consequence; the residual $\rho_\env - (1-\alpha)\rho_\tot$, which is the
parameter draw's price in environment value units and is the missing piece of \S\ref{sec:meas-g};
and $\Cov_{\pi}(Q_\tot, Q_\env)$. It is instrumentation, it is not required by \eqref{eq:rule},
and it should be built when one of those three is wanted rather than on this document's account:
\begin{equation}
  \hat{Q}_\env(s_t,a_t) = R^{(n)}_t + \gamma^{n} \ip{\pi}{Q_{\env,\theta'}}(s_{t+n}).
  \label{eq:qenv}
\end{equation}

\section{Extensions readmitted by the last two cases}
\label{sec:readmit}

\S\ref{sec:c10} and \S\ref{sec:c11} are failed by everything that runs for free, and the second
of them is the whole late half of a run. The previous revision excluded a head, a servo and a
normaliser by price. That was the wrong test --- the right one is whether a case demands the
thing --- and applied to these two cases it lets each back in for a reason the earlier version
did not have. None of what follows is required by \eqref{eq:rule}, which is why it is a separate
section and why \S\ref{sec:staging} does not wait on it. One condition applies to all of it:
a parameter can only be checked through its own output, so whatever one is meant to measure gets
logged beside it.

\subsection{Fixing the parameter draw}
\label{sec:fixdraw}

Two changes follow from Remark~\ref{rem:peraction} and neither needs a parameter. Give
$\varepsilon_{\mathrm{in}}$ and $\varepsilon_{\mathrm{out}}$ a batch axis, so that the
environments explore independently instead of being an ensemble of one; and derive the draw from
a key folded on a per-environment episode counter, so that it is held while its consequences play
out. The first is a per-sample product where there is now a broadcast. The second needs the
terminal flags the loop already handles.

This was the cheapest item before the reading and it is the gating item after it, because three
separate things turn out to rest on the draws being posterior samples: \S\ref{sec:c10}'s split
and \S\ref{sec:c11}'s contracting signal, both through \eqref{eq:uasplit}, and the claim ---
which Dabney et al.\ make on the whole NoisyNets family (\S\ref{sec:related}) --- that this kind
of exploration is temporally persistent at all. If the draw is fixed, the two cases that nothing
free passes become affordable at one extra output-layer pass. If it is not worth fixing, then
dropping \texttt{NoisyLinear} returns a large fraction of the parameters and makes the behaviour
policy the
only exploration mechanism in the loop, which is the condition under which $\rho$ becomes
directly measurable instead of resting on a floor that is most of itself
(\S\ref{sec:meas-rho}).

\emph{And there is a third option that skips the question.} $\epsilon z$-greedy buys temporal
persistence with no network change at all, by repeating the exploratory action for a
heavy-tailed duration (\S\ref{sec:related}). It is cheaper than either branch above and it is
the control that a persistence experiment should be run against, since if repeating a uniform
action for $\zeta$-distributed durations recovers whatever a held parameter draw buys, the
parameter draw is not what was needed.

\subsection{The cost of a contracting signal}
\label{sec:contracting}

\S\ref{sec:c11} asks for something that goes as $1/\sqrt{n}$. In descending order of cheapness,
and the first is new since the reading:

\emph{(0) Two draws of the existing head}, through \eqref{eq:uasplit} --- \textbf{measured, and
it does not work as things stand}. One extra pass through the output projections returns the
epistemic and aleatoric halves separately and without bias, which is \S\ref{sec:c10} and
\S\ref{sec:c11} at once, and it costs no parameters. \S\ref{sec:denoms} ran it: the estimator
behaves (its covariance term is positive everywhere, and it agrees with the naive one on the
aleatoric half closely), and the epistemic half it recovers is flat across a run and fails to
transfer across checkpoints (\S\ref{sec:denoms}). So the failure is located: it is the draws, not
the
readout. This is the cheapest thing on the list and it becomes the whole of \S\ref{sec:c11}'s
answer the moment \S\ref{sec:fixdraw} lands --- and only then.

\emph{(a) A last-layer posterior.} Treat the $512$-wide activation as a fixed basis and keep a
Gaussian posterior over the head weights, updated in closed form; act on a draw held for an
interval. Contracts as $1/\sqrt{n}$ by construction, per-action, no new network --- a covariance
per action is megabytes at this width. It also supplies the draws (0) needs. The head is
dueling and distribution-valued, so ``the feature'' needs defining, the activation averaged over
the head's own risk axis being the obvious choice.

\emph{(b) Random-prior bootstrap heads.} Parameters are not the objection --- a final projection
per member is negligible. The objections are that members sharing one feature layer have
little left to disagree about, which is what the random prior exists to fix, and that they would
be output layers in a network whose unfixed pathology is output-layer effective-learning-rate
collapse. \S\ref{sec:related} shows this is the configuration deep IDS runs on Atari, down to the
\texttt{stop\_gradient} on the distributional branch.

\emph{(c) An intrinsic-value head.} A novelty signal --- RND's prediction error against a fixed
random target, or a latent-prediction error of the BYOL-Explore kind, which this codebase is
closer to having since it already carries a self-predictive branch on the plan of record ---
carried through its own value head gives a per-\emph{state} quantity that contracts as a state
becomes familiar. It is the only one that says anything about states the agent has not reached,
and it is the only one whose signal is not in value units, so it is the one that needs a carrier
(\S\ref{sec:c1}'s $\ddagger$). It belongs to \texttt{btr.py}'s exploration backlog rather than to
this document, except through \S\ref{sec:coupling}.

\emph{The shape to put any of them in} is the one \S\ref{sec:denoms} forced by separating what
$\dbar$ was doing badly into two jobs: a carrier with the units, and a dimensionless modulation
with the epistemics. Boltzmann--Gumbel's $\sigma/\sqrt{N_a}$ --- the reward's subgaussian
parameter over an action's own visit count --- is that shape with a regret bound behind it, and
Fox's $2\mu_A/\sigma_A^2$ --- the mean of the gap's posterior over its variance --- is that shape
with a bias argument behind it (both \S\ref{sec:related}); they disagree about the exponent,
which \S\ref{sec:owed} now carries.

\subsection{Coupling under a contracting denominator}
\label{sec:coupling}

\S\ref{sec:c8-staging} refuses to couple \eqref{eq:rule} because the sharp limit has to stay
reachable, which is a fact about staging and not about the denominator. With (I)'s denominator
the same mechanism is the thing one would pay a head for, so it is worth stating as an experiment
rather than leaving as a remark --- with the ceiling that \S\ref{sec:meas-couple} now puts on it.

\emph{The arm.} $\tau(s) = \hat{s}(s)/g$ in \texttt{\_act} \emph{and} in the Munchausen target,
$\hat{s}$ contracting, against the same rule decoupled. The two arms differ in exactly one term
--- \eqref{eq:distortion} evaluated at $\hat{s}(s)/g$ --- so the comparison measures the intrinsic
reward and nothing else, which is what an intrinsic-motivation scheme should get: a deliberate
experiment with its own baseline.

\emph{What it would buy.} An exploration bonus that is discounted, propagated \emph{and spent}:
the value of a state that \emph{leads} to unresolved states rises, and falls again once they are
resolved. Coupling any denominator buys the first two (\S\ref{sec:c8}); only a contracting one
buys the third, and without it the same machinery yields a permanent pull toward states nothing
can settle. It costs
no intrinsic-reward head beyond whatever supplies $\hat{s}$, no second value stream, and it adds
nothing to any individual action's value --- the denominator is a per-state scale, so the pitfall
of optimism (\S\ref{sec:related}) does not apply to it.

\emph{What it cannot buy, which is new and is measurable before the run.} The bonus reaches the
objective only through \eqref{eq:distortion}, which is exponentially gated by the state's own
gaps: hold every state at the ceiling and the great majority still receive under a tenth of
$\tau_\ell\ln|\actions|$ (\S\ref{sec:meas-couple}). A state that is genuinely unvisited
but where the critic is confidently wrong is paid nothing, which is precisely the state a novelty
bonus exists to reach, and which a reward-level bonus reaches because it is ungated.
\textbf{So the first thing to do is not the run.} It is to take whatever $\hat{s}$
\S\ref{sec:fixdraw} produces, evaluate \eqref{eq:distortion} on it over replay states with
\texttt{scripts/coupling\_cost.py}, and see whether the bonus lands on the states it is meant for
or on the near-ties, as it does for \eqref{eq:rule}. That costs minutes and it decides whether
the arm is worth a baseline.

\emph{What it costs, and what has to be checked first.} The control path, which is the expensive
part: a coupled arm needs its own baseline and forfeits the full-budget references
(\S\ref{sec:project}). Three things then need re-examining, none of them fatal and none of them
free. $l_0 = -1$ is \emph{not} one of them, which \S\ref{sec:couple-case} settles: coupled, a
drawn transition is clipped with probability at most $|\actions|e^{-1/(\alpha\tau(s))}$ at each
state, so the clip stays inert for any $\hat{s}/g$ sharp enough to be worth running and wakes
only for a rule an order more diffuse than \eqref{eq:rule}.
\textbf{The fixed point moves with $\lambda(s)$, and once
$\lambda$ is a functional of $q$ the operator is not the one
Remark~\ref{rem:coupled-ok} analyses} --- that remark licenses a fixed per-state schedule and no
more, so a coupled arm is running an operator with no contraction argument behind it and should
log $\srank$'s own trajectory to catch the feedback loop if it runs away. And $\hat{s}$ is then
inside the loss, so a bug in it is a bug in the objective rather than a bug in the sampling,
which is a different class of failure to debug. \S\ref{sec:fixdraw} comes first regardless: there is no
contracting $\hat{s}$ on this codebase until it does.

\subsection{Servo control, where it is well posed}
\label{sec:servo}

What retired the three servos stands exactly where it was aimed: at servoing the \emph{level}
against $J$, which is not monotone and is worth less than a seed. Neither objection touches a
servo on a monotone, cheaply measured quantity, and Proposition~\ref{prop:rho} is precisely that
certificate for regret against the scale. So servo $g$ --- not the return --- against a
behavioural target measurable at every step: an off-greedy rate defined on values rather than on
the argmax, so that \S\ref{sec:c3} survives, or --- once the draw has been dealt with and its
floor no longer swamps the measurement --- the realised $\rho/\gapbar$. One integrator, one
constant, no head, and it converts the rule's last hand-set number into a target.

The measured warning is new and it inverts the previous revision's assumption. Rates were
supposed to be the portable anchor and value quantities the unportable ones. In fact the two best
arms this project has run sit an order of magnitude apart in off-greedy rate
(\S\ref{sec:meas-g}), so \emph{the target does not
transfer between games} even though \eqref{eq:rule} does. A servo therefore sets the level within
a game; what crosses games is the rule, not its set point.

\subsection{Return normalisation}
\label{sec:norm}

True of \eqref{eq:rule}'s constants, false of everything around them. A running estimate
$\snorm$ of the dispersion of returns in the buffer costs one scalar and one update; three
things it buys.

\emph{It is what the primed candidates are made of, and the table says how far that gets.}
$\snorm$ is the only thing that repairs (C)'s \S\ref{sec:c1} failure, and (C) is the most
precisely right candidate on \S\ref{sec:c2}: the mistake gets rarer at exactly the rate that
keeps its cost constant. So (C$'$) is carried through \S\ref{sec:cases} rather than introduced
here, and what the column shows is that the repair is real and partial --- \S\ref{sec:c1} is
fixed, while \S\ref{sec:c5} and \S\ref{sec:c6} are not, a budget fixed in units of the game's
return dispersion being still infeasible once the gaps at a state contract under it. A global
normaliser divides out the value unit once per game; it cannot divide it out once per state,
which is the difference between it and \eqref{eq:rule}.

\emph{It makes the inherited absolute constants mean something.} \texttt{IQN\_HUBER\_LOSS\_K} $=
1$ and $l_0 = -1$ are both thresholds inherited from reward clipping's $[-1,1]$. Not this
document's problem, but the same normaliser is what fixes them, and the plan of record already
schedules one for the categorical arm.

\emph{And only an external scale can detect a collapse.} Every ratio in \eqref{eq:rule} is
homogeneous of degree zero, which is exactly why it passes \S\ref{sec:c6} --- and exactly why it
is blind there. $\gapbar/\srank$ cannot distinguish a fine-grained game from a value function
that has died, because both look like small gaps against a matching spread. $\snorm$ is
estimated from returns in the buffer rather than from the network, so $\gapbar/\snorm$ can.
\eqref{eq:rule} declines to amplify the Enduro failure; an external unit is what it would take to
push back on it, and the diagnostic is worth having even if nothing acts on it.

\emph{And it is nearly the same object as \S\ref{sec:denoms}'s}, which makes both cheaper to
justify: $\sret$ is a per-state PopArt scale obtained for free, and a global $\snorm$ is its
average over states. Normalising globally and acting at a fixed $\tau$ in normalised units is
candidate (A$'$), which is this document's rule at coarser resolution. The per-state refinement
is measured to be worth about what the cross-game one is (\S\ref{sec:denoms}) --- so they are
complements and not
alternatives.

\subsection{The assembled scheme}

Information-directed sampling, whose deep form \S\ref{sec:related} states in full. Three
properties earn it the place. It acts at the argmin over actions of a ratio $\hat\Delta^2/I$,
a squared per-action regret estimate over an information gain (\S\ref{sec:related} states both;
neither symbol is this document's $\rho$ or $\lambda$). It is scale-invariant by construction,
since rescaling rewards by $c$ multiplies that ratio by $c^2$ uniformly and the argmin does not
move; it prices a catastrophic action by the square of its regret and so passes \S\ref{sec:c2}
with no floor; and it is the only
candidate in this document that trades regret against information rather than assuming a
functional form for the trade. Its two inputs are the two halves \S\ref{sec:c10} asks to be
separated, and \eqref{eq:uasplit} is a cheaper way to get them than the one it uses. It is named
as the destination rather than proposed: the components land one at a time, with \eqref{eq:rule}
as the control.

\section{Related work}
\label{sec:related}

Ten papers, read rather than searched. Three of them change a claim made above and are marked
\textbf{[changes]}; the rest place the design.

\paragraph{\textbf{[changes \S\ref{sec:denoms}, \S\ref{sec:rule}]} The rule's functional form is
not new, and the exponent is now the open question.} Fox, \emph{Toward provably unbiased
temporal-difference value estimation} (NeurIPS 2019 workshop), reported in Hu, Abbeel and Fox
(\texttt{arXiv:2111.14204}): in a two-action problem where the gap $Q(s,a_0) - Q(s,a_1)$ has
Gaussian uncertainty with mean $\mu_A$ and variance $\sigma_A^2$, the inverse temperature that
eliminates the bias in an entropy-regularised value update is, in closed form,
$\beta(s) = 2\mu_A/\sigma_A^2$. A per-state temperature set by a gap over the spread of that gap
therefore has a derivation behind it and predates this document. Two differences survive, and
both are now owed rather than settled:

\begin{itemize}
  \item \emph{Which $\sigma$.} Fox's is the epistemic variance of the gap \emph{estimate}, and
  \S\ref{sec:denoms} measures $\srank$ to be essentially aleatoric. So the two rules are not
  variants of each other; they divide by different quantities that happen to share a shape.
  \item \emph{The exponent, which only matters once the first point is settled.} In the
  sharpness units of \S\ref{sec:bg}, where \eqref{eq:rule}'s operating branch reads
  $k = g\,\gapbar/\srank$ (\eqref{eq:rule-k}), Fox's rule is $2(\mu_A/\sigma_A)^2$ --- the
  \emph{square} of the signal-to-noise ratio, where the rule is linear in it. Substituted with the
  epistemic spread this codebase actually has, it fails to transfer by a wide margin where
  $\srank$ transfers (\S\ref{sec:denoms}), because squaring a denominator that does not
  transfer does not transfer twice as badly but four times. That is not a verdict on Fox's
  exponent --- it is the same verdict on the same draw, arrived at through a third statistic.
\end{itemize}

The same note supplies what \S\ref{sec:c11} asks for and no statistic in \S\ref{sec:denoms}
gives: by the central limit theorem $\sigma_A^2$ falls as $1/n$, so Fox's $\beta$ grows linearly
in the evidence at a state. Hu et al.\ then abandon the closed form for exactly that asymptotic
--- their CBSQL sets $\beta(s) = \kappa\, n(s)$ from a CTS pseudo-count --- and, worth noting for
\S\ref{sec:decouple}, they put it in the \emph{backup} (a mellowmax target) and act
$\epsilon$-greedy. The one published per-state temperature schedule is in the loss, not in the
behaviour policy.

\paragraph{A state-dependent Boltzmann temperature by root-find also exists, and it is degenerate
exactly where \S\ref{sec:c9} says.} Asadi and Littman, \emph{An alternative softmax operator}
(ICML 2017, \texttt{arXiv:1612.05628}). Their maximum-entropy mellowmax policy is
$\sm(\beta \hat{Q}(s,\cdot))$ with $\beta$ the root of
$\sum_a e^{\beta(\hat{Q}_a - \mmax_\omega \hat{Q})}(\hat{Q}_a - \mmax_\omega \hat{Q}) = 0$, solved per
state by Brent's method. At a tied state every bracket vanishes and the equation reads $0 = 0$
for every $\beta$ --- the degeneracy \S\ref{sec:c9} identified in the previous revision's
root-find, in the published version of the same construction. Their motivation is orthogonal to
this document's: $\beta(s)$ is chosen so the induced operator is a non-expansion and GVI has a
unique fixed point, not so the constant transfers. They also show the fixed-$\beta$ Boltzmann
operator has \emph{multiple} fixed points on a two-state MDP, which is a hazard
Remark~\ref{rem:coupled-ok}'s contraction result does not cover --- and it is the second reason
that remark cannot license a $\beta(s)$ read off $q$, the first being that it assumes one that is
not.

\paragraph{\textbf{[changes \S\ref{sec:c10}, \S\ref{sec:readmit}]} The epistemic/aleatoric split
is purchasable for two draws, not for a head.} Clements et al., \emph{Estimating risk and
uncertainty in deep RL} (\texttt{arXiv:1905.09638}). For a quantile network with outputs
$y_i(\theta)$ they define $\sigma^2_{\mathrm{epi}} = \E_i \Var_\theta(y_i)$ and
$\sigma^2_{\mathrm{ale}} = \Var_i(\E_\theta y_i)$, prove
$\Var_{\theta,i}(y_i) = \sigma^2_{\mathrm{epi}} + \sigma^2_{\mathrm{ale}}$, and give the unbiased
two-draw estimators of both that \S\ref{sec:bg} states as \eqref{eq:uasplit},
whose own variance vanishes as the number of quantiles grows. Two draws is what the previous
revision was already willing to pay for $\dbar$, and the acting quantile sample is what
\texttt{\_act}
already has. They also prove the obvious single-draw estimator --- $\Var_i[y_i(\hat\theta)]$,
which is $\sret$ --- is biased \emph{upward} as an estimate of the aleatoric part, by the
epistemic term. \S\ref{sec:denoms} now runs all of this. The estimator behaves --- the covariance term is
positive at every state, and the aleatoric half agrees closely with the same spread taken on
the noise-zeroed network, so the bias they warn about does not reach the figures here. It does
not because of the zeroing and not because the epistemic term is small --- on the centred
statistic it is not small (\S\ref{sec:denoms}) --- and the epistemic term it recovers is flat
across a run and spans
widely across games and behaviour policies. \eqref{eq:uasplit} is only as good as the
draws, and it inherits Remark~\ref{rem:peraction} entire; their own draws come from two
auxiliary networks under randomised MAP sampling, which is the cost they pay and we have not.
The one part of the construction that is quantile-specific --- the pairing at a common quantile
index --- is therefore also the part worth nothing here, which is convenient given that
\S\ref{sec:bg} defines $\srank$ so that it survives a categorical head.

\paragraph{\textbf{[changes \S\ref{sec:meas-g}]} Churn is a third exploration source of the same
size as the one being calibrated.} Schaul, Barreto, Quan and Ostrovski, \emph{The phenomenon of
policy churn} (\texttt{arXiv:2206.00730}). In DQN on Atari the greedy action changes in
${\sim}10\%$ of states per \emph{single} gradient update, which over a run is a million changes
per state; they argue this implicit exploration is why $\epsilon$ matters less than assumed. BTR's
Table~5 measures $3.8\%$ churn on Phoenix with IQN, against the off-greedy rate
\S\ref{sec:meas-g} anchors $g$ on. So the anchor is being matched to a quantity of the same
magnitude as one nobody controls, which is a caveat on the anchor and not on the rule. One
further reading cuts against \S\ref{sec:c3}: they test churn with the full Atari action set
exposed, creating deliberate redundancy, and find the magnitude barely moves, with no structure
relating swaps to known action equivalences. Duplication may therefore explain less of Zaxxon's
small gap than \S\ref{sec:c3} assumes.

\paragraph{A per-action rate from a reward scale over a count, with a regret bound --- in
bandits.} Cesa-Bianchi, Gentile, Lugosi and Neu, \emph{Boltzmann exploration done right}
(\texttt{arXiv:1705.10257}). The negative result is sharper than the abstract suggests and also
narrower: in a two-armed \emph{stochastic bandit}, Boltzmann exploration with any learning-rate
sequence growing faster than $2\log t$ has linear regret, and monotone schedules are suboptimal
generally. Their remedy, $\beta_{t,i} = \sqrt{C^2/N_{t,i}}$ with $C$ the reward's subgaussian
parameter, is per-\emph{arm}, and it is not a Boltzmann policy: they perturb each estimate by an
independent Gumbel and take the argmax, which coincides with a softmax only when the scale is
common across actions. So the carrier-times-modulation shape is right and the softmax is not the
form it takes. Two details transfer: the bound survives underestimating $C$ (penalised
exponentially in the true $\sigma^2$), and a heavy-tailed variant replaces $\sigma^2$ by a bound
on the second moment, so the scale need not be known exactly. None of it is a discounted control
result, which is what \S\ref{sec:owed} wanted to know.

\paragraph{The quantile spread is a mixture, and using it naively is a known failure.} Mavrin
et al., \emph{Distributional RL for efficient exploration} (\texttt{arXiv:1905.06125}). A QR
distribution mixes parametric uncertainty, which vanishes with data, and intrinsic uncertainty,
which does not; they construct a bandit where a bonus from the raw spread never selects the best
arm, because the best arm has the lowest intrinsic variance. Their two patches are a decaying
multiplier $c_t = c\sqrt{\log t/t}$ --- the rate at which parametric uncertainty decays in
quantile-regression theory, and a global clock of exactly the kind Cesa-Bianchi et al.\ prove
suboptimal --- and a left-truncated variance about the median, for asymmetry. \S\ref{sec:denoms}
is their prediction measured: \S\ref{sec:denoms} finds $\sret$ flat across checkpoints, which
is
what ``the parametric part has already vanished and the intrinsic part remains'' looks like.

\emph{Why their counterexample does not transfer, and what does.} Their bonus is per-action and
\emph{added to the value}, so it biases which action is chosen toward the noisy one; Cho et al.,
\emph{Pitfall of optimism} (\texttt{arXiv:2310.16546}), make the same point as their central
claim and fix it by randomising the risk criterion instead. \eqref{eq:rule} adds nothing to any
action's value: its denominator is a mean over actions, a per-state scale, and a scale cannot
change the ordering. What survives of the objection is weaker and is \S\ref{sec:c10}'s version
--- the policy diffuses where outcomes are noisy, rather than preferring the noisy action.

\paragraph{The assembly, with the normaliser this document arrived at independently.} Nikolov,
Kirschner, Berkenkamp and Krause (\texttt{arXiv:1812.07544}). Their deep IDS is, per step,
$\hat\Delta(s,a) = \max_{a'}[\mu + \lambda\sigma](s,a') - [\mu - \lambda\sigma](s,a)$,
$I(s,a) = \log(1 + \sigma^2(s,a)/\rho^2(s,a)) + \epsilon_2$, act at $\argmin_a \hat\Delta^2/I$,
with $\sigma$ from $K$ bootstrap heads split after the convolutions and $\rho^2$ from a separate
distributional head off the shared features, its loss not propagated into the trunk. Two things
to take. Their $\rho^2(s,a)$ is $\Var(Z(s,a))$ \emph{divided by the mean over actions of
$\Var(Z(s,\cdot))$} --- which is $\sret$, used as a per-state scale to make the ratio
dimensionless, arrived at here from the units and there from the information ratio. And the
architecture is exactly \S\ref{sec:readmit}(b): bootstrap heads on the trunk, distributional head
behind a \texttt{stop\_gradient}.

\paragraph{Persistence is buyable without touching the network at all.} Dabney, Ostrovski and
Barreto, \emph{Temporally-extended $\epsilon$-greedy} (\texttt{arXiv:2006.01782}). Repeat the
exploratory action for a duration $n \sim \zeta(\mu)$, heavy-tailed, $\mu = 2$, capped at $100$
for Rainbow-scale agents; comparable to pseudo-counts, RND and bootstrap on hard-exploration
Atari without their losses elsewhere. It is carried in \S\ref{sec:fixdraw} as the control that any persistence experiment here should
be run against. Note also that they classify NoisyNets \emph{with}
Thompson sampling and randomised priors as ``non-stationary, temporally persistent'' --- the
property Remark~\ref{rem:peraction} finds this codebase's wiring does not have. The family is
credited with persistence it is not getting here.

\paragraph{And what the search did not find, restated.} No paper sets a Boltzmann temperature
from a \emph{distributional} spread per state. Fox has the shape from the epistemic variance of
the gap; Mavrin has the distributional spread as a bonus; Asadi and Littman have the per-state
root-find from an operator; Nikolov has the mean-over-actions return variance as a normaliser
inside an information ratio. \eqref{eq:rule} is the one combination none of them writes down,
and after reading, the reason looks less like an oversight and more like the fact that everyone
who examined the spread closely concluded it was the wrong half of a mixture. \S\ref{sec:denoms}
agrees with them about which half $\srank$ is. What it adds is that on this codebase the other
half is not available at any price short of \S\ref{sec:fixdraw}, and that the wrong half is the
only one that transfers between two games.

\section{Measurements}
\label{sec:meas}

Everything measured in this document is here, and nothing above is measured anywhere else. There
are four sources. Per-state statistics read off four checkpoints with the noise zeroed
(\texttt{scripts/spread\_probe.py}, \texttt{scripts/probe\_checkpoint.py},
\texttt{results/spread\_*.json}, \texttt{results/probe\_*.json}, \texttt{results/rho\_*.json}):
minutes on a checkpoint, reading the network rather than the return. Policies scored by rollout
at fixed weights (\texttt{scripts/score\_policies.py}): hours, and noisy, as \S\ref{sec:meas-return}
says. Post-processes of the per-state arrays the first of those writes
(\texttt{scripts/coupling\_cost.py}, \texttt{results/coupling\_cost.json}): seconds, no
checkpoint, and re-runnable against any rule stated as a $\tau(s)$. And four complete $200$M-frame
training runs, which are the expensive ones and the reason \S\ref{sec:obj}'s criterion is set
where it is.

One caveat governs all of it and is not repeated below. \emph{The four checkpoints are four
different runs} --- two games, two behaviour policies, two points in training --- so every
cross-checkpoint statement here is between runs and not within one. Whether the ratio this
document is built on is stable inside a single run has never been watched, and it is
\S\ref{sec:staging}'s first stage precisely because it is free.


\subsection{The setting's own movement, and two verifications}
\label{sec:meas-mvi}

\paragraph{What a constant in value units would have to survive.} Over a single run
\texttt{train/q} moves $195\times$ and \texttt{train/action\_gap} moves $13\times$, and their
ratio $\mathrm{gap}/q$ \emph{falls} $7\times$ rather than holding. Across games the top-two gap is
$0.033$ on Zaxxon against $0.179$--$0.231$ on Phoenix. This is \S\ref{sec:project}'s second
constraint measured, and it is why no rule in this document carries a constant in value units.

\paragraph{Proposition~\ref{prop:amp}, numerically.} \texttt{scripts/check\_mvi\_gap.py} iterates
M-VI and entropy-regularised VI to convergence on random finite MDPs. At
$\alpha \in \{0.5, 0.9, 0.99\}$ the measured gap ratio is $2$, $10$ and $100$ to six decimals at
every state-action pair, against Munchausen's published $3$, $19$ and $199$
(Remark~\ref{rem:erratum}), and $\|\sm(q_\infty/\tau) - \pi_*^\lambda\|_\infty < 10^{-13}$. The
same script confirms Remark~\ref{rem:coupled-ok}'s contraction with $\tau(s)$ drawn once and
spanning a decade.

\subsection{The observed value collapse}
\label{sec:meas-collapse}

The event \S\ref{sec:c6} is built on, and the only one of this document's three observed cases
that is a single incident rather than a trend. On the $\texttt{ACT\_TEMPERATURE} = 0.01$ Phoenix
arm, \texttt{train/q} fell from $3.24$ to $1.97$ between $5$M and $11$M env steps with the policy
going near-uniform, and \texttt{eval/returns} stayed flat to $33$M frames --- some $20$M frames
given back, which is most of what sharpening from $0.03$ had bought (\S\ref{sec:meas-return}).
The softmax's self-correction only reopens once the gap falls to $O(\tau)$, so a sharper draw
carries a deeper collapse: this is the mechanism \S\ref{sec:c6} abstracts and the reason
\eqref{eq:rule} needs the floor branch at all.

\subsection{Candidate denominators}
\label{sec:denoms}

Every case in \S\ref{sec:cases} asks for a scale that comes from the state, and \eqref{eq:rule}
names one. This subsection is what was available to supply it and why that one was chosen. Four checkpoints, $2048$ states from each one's own rollout, per-state statistics
with the noise zeroed and $\dbar$ over eight live draws; \texttt{scripts/spread\_probe.py} and
\texttt{results/spread\_*.json}, minutes on a checkpoint, reading the network rather than the
return.

Every candidate is a mean over actions, so \S\ref{sec:c3} holds by construction. $\dbar$ is the
previous revision's, the spread of $\tilde{q}$ across draws of the parameter noise. $\sret$ and
$\srank$ are read off the return distribution the head already emits. $\tilde\sigma_{\mathrm{epi}}$
and $\tilde\sigma_{\mathrm{ale}}$ are \eqref{eq:uasplit} over the same eight draws, on the raw
distribution and on the centred advantage alike.

\emph{One correction to make first.} The draws must share a quantile draw. Estimating $\dbar$
with a fresh $u$ sample per pass charges quantile-sampling noise to the parameter draw and
inflates it by some $75\%$; every figure below fixes the sample across draws, and $\dbar$ is
correspondingly smaller than an earlier reading of it.

\begin{center}
\small
\begin{tabular}{lcccc}
  \toprule
   & Phoenix $46$M & Phoenix $250$M & Phoenix $200$M & Zaxxon $204$M \\
   & soft $0.03$ & soft $0.005$ & argmax & soft $0.03$ \\
  \midrule
  $\gapbar$              & $0.405$  & $0.721$  & $0.934$  & $0.160$ \\
  top-two gap            & $0.104$  & $0.187$  & $0.204$  & $0.033$ \\
  $\sret$                & $0.439$  & $0.433$  & $0.432$  & $0.206$ \\
  $\srank$               & $0.0366$ & $0.0736$ & $0.0651$ & $0.0171$ \\
  $\dbar$                & $0.0540$ & $0.0526$ & $0.0590$ & $0.0282$ \\
  \midrule
  $\tilde\sigma_{\mathrm{epi}}$      & $0.0743$ & $0.0661$ & $0.0729$ & $0.0345$ \\
  $\tilde\sigma_{\mathrm{ale}}$      & $0.4366$ & $0.4324$ & $0.4327$ & $0.2064$ \\
  $\tilde\sigma_{\mathrm{ale}}$, centred & $0.0363$ & $0.0730$ & $0.0648$ & $0.0169$ \\
  $\tilde\sigma_{\mathrm{epi}}$, centred & $0.0250$ & $0.0349$ & $0.0339$ & $0.0150$ \\
  \midrule
  $\gapbar/\sret$        & $0.99$   & $1.77$   & $2.35$   & $0.83$ \\
  $\gapbar/\srank$       & $11.78$  & $11.24$  & $16.54$  & $11.21$ \\
  $\gapbar/\dbar$        & $7.78$   & $14.35$  & $16.76$  & $5.77$ \\
  $\gapbar/\tilde\sigma_{\mathrm{epi}}$ & $5.48$ & $11.18$ & $13.18$ & $4.61$ \\
  $2(\gapbar/\tilde\sigma_{\mathrm{epi}})^2$ & $69$ & $261$ & $387$ & $50$ \\
  $\eta^2$               & $0.315$  & $0.542$  & $0.657$  & $0.164$ \\
  \bottomrule
\end{tabular}
\end{center}

Means over states; the ratios are means of the per-state ratio, which is what a per-state rule
divides by, and not ratios of the columns above them. $\eta^2(s)$ is the share of the return's
variance at $s$ that the choice of action explains, and is in the table as the one line that says
how much of what the critic emits is about the decision at all.

\paragraph{The prerequisite fails, on both readings it was given.} $\dbar$ does not move with the
value scale: $0.0540$, $0.0526$, $0.0590$ across three Phoenix checkpoints over which $\gapbar$
moves $2.3\times$. It does not fall as a run progresses. And across games it moves $1.9\times$
where $\gapbar$ moves $4.5\times$, so $\gapbar/\dbar$ spans $5.8$ to $16.8$ over the four and a
gain calibrated on one is wrong by up to $2.9\times$ on another. That is \S\ref{sec:project}'s
second constraint failed by the denominator itself, which is the one place this design cannot
absorb it.

\paragraph{And the properly defined epistemic spread fails the same way, which closes the
question.} $\dbar$ is a crude statistic --- the spread of the \emph{mean} across draws --- so
its failure could have been the estimator rather than the quantity. \eqref{eq:uasplit} is the
unbiased one, and it agrees: $\tilde\sigma_{\mathrm{epi}}$ reads $0.0743$, $0.0661$, $0.0729$ on
the three Phoenix checkpoints, flat again, and $\gapbar/\tilde\sigma_{\mathrm{epi}}$ spans
$4.6$--$13.2$ over the four, which is no better than $\dbar$'s and far worse than $\srank$'s
$1.5$.
Fox's form fares worse still, its denominator being squared: $2(\gapbar/\tilde\sigma_{\mathrm{epi}})^2$
spans $50$ to $387$ (\S\ref{sec:related}). \textbf{No epistemic denominator available on this
codebase transfers, and the estimator is not what is wrong with them.} The draw is
(Remark~\ref{rem:peraction}).

\paragraph{$\srank$ is measured to be aleatoric, which confirms \S\ref{sec:c10} rather than
answering it.} Taking \eqref{eq:uasplit} on the centred advantage, the aleatoric half reads
$0.0363$, $0.0730$, $0.0648$, $0.0169$ against $\srank$'s $0.0366$, $0.0736$, $0.0651$, $0.0171$
--- equal to within $1\%$ on every checkpoint. The epistemic half of the ordering's spread is
half to nine-tenths of that in magnitude and is the part that does not transfer. The same holds
for $\sret$ against $\tilde\sigma_{\mathrm{ale}}$, where the epistemic term is $15$--$17\%$ of
the aleatoric one and is genuinely small. Two consequences. \emph{Zeroing the noise is what buys
the agreement}, not the epistemic term being small: the bias \S\ref{sec:related} warns of is the
bias of a spread taken \emph{on a draw}, and every figure here is taken with the noise zeroed,
which is why $\srank$ lands on the aleatoric half instead of above it --- on a live draw the
centred statistic would sit a tenth to a third high. And the decomposition survives the centring,
the covariance estimator being positive at every state on every checkpoint.

\paragraph{And it is flat for a reason that is worse than drift.} \texttt{noisy\_sigma\_ratio}
falls from ${\sim}0.6$ to $0.024$--$0.071$ over a run while \texttt{weights/global\_norm} goes
$61 \to 162$; the product of a decaying $\snoise$ and a growing kernel is roughly constant.
$\dbar$ is therefore not an uncertainty that happens to be flat, it is two unrelated drifts
cancelling --- and the project's stated next intervention is on the second of them. A rule
divided by $\dbar$ would have had its exploration schedule silently rewritten by a fix aimed at
\texttt{adam\_step\_rel}.

\paragraph{$\sret$ is a property of the game.} $0.439$, $0.433$, $0.432$ on Phoenix regardless of
epoch or behaviour policy, against $0.206$ on Zaxxon. It is the return's own dispersion and it is
the right \emph{unit} --- a per-state PopArt scale, obtained for free (\S\ref{sec:readmit}) ---
but a temperature $\sret/g$ is then a per-game constant that calibrates itself, not a rule that
responds to a state. Its per-state variation is real ($0.96$--$2.95$ in $\gapbar/\sret$ at
$\text{p}5$--$\text{p}95$ on the $250$M Phoenix net) but \S\ref{sec:c10} says the variation is the wrong kind.

\paragraph{$\srank$ is a property of how well the ordering is resolved, and it transfers.}
$\gapbar/\srank$ reads $11.78$, $11.24$ and $11.21$ on three soft-trained checkpoints spanning
two games, $4.5\times$ in $\gapbar$, $8$ actions against $18$, and $46$M frames against $250$M.
The whole per-state distribution coincides, not only the mean: $\text{p}5$--$\text{p}95$ is
$4.5$--$18.8$ on Phoenix against $4.7$--$20.3$ on Zaxxon. It moves to $16.54$ on the
argmax-trained net, which is the arm that scores $1.30\times$ the soft one and whose ordering is
by that much better resolved --- so the quantity is not constant, it tracks the one thing it is
supposed to track and is invariant to the things it should be invariant to. Two games is two
games, and this is the claim most worth breaking on a third.

\paragraph{Both distributional denominators carry information the gap does not.} On the two
Phoenix nets the correlation between $\ln\gapbar$ and $\ln\srank$ over states is $+0.20$ and
$+0.06$, so the ratio is not a restatement of the numerator; on Zaxxon it is $+0.63$ and half of
it is. $\dbar$ is comparably independent ($+0.06$, $+0.05$, $+0.26$) and correlates
$0.47$--$0.67$ with the distributional ratio, so the two denominators agree about half the time and disagree
enough that the choice between them is not cosmetic.

\paragraph{Nothing here anneals.} This is \S\ref{sec:c11}, and it is the finding that decides
what the rest of the document is about. $\gapbar/\srank$ is flat from $46$M to $250$M, $\dbar$ is
flat in absolute terms, $\sret$ is flat in absolute terms. A rule built from any of them holds a
fixed \emph{relative} sharpness for an entire run. That is a strict improvement on holding a
fixed absolute one and it is not the same thing as exploring less as you learn more.


\subsection{The distortion family}
\label{sec:meas-family}

\S\ref{sec:bg} demands translation equivariance and positive homogeneity of every member of the
family $\srankb$ is taken over. This is those demands checked, and what the choice costs once
they are met. \texttt{scripts/distortion\_family.py}: seconds, no checkpoint and no rollout, every
figure on constructed distributions whose answer is known in closed form --- which is the point,
since an inadmissible reading is inadmissible for every network and there is nothing to learn
from a checkpoint about it.

\paragraph{Equivariance, on a pair whose answer is zero.} Two actions whose return distributions
are exact translates: the ordering is resolved at every risk level, so $\srank = 0$ by
\S\ref{sec:bg}'s zero set. Read through a categorical head:

\begin{center}
\small
\begin{tabular}{lcc}
  \toprule
  reading & $\srankb$ & \\
  \midrule
  atom probabilities $p_i$          & $0.0246$ & inadmissible \\
  tail means, threshold fixed on the support & $0.0563$ & inadmissible \\
  tail means, each action's own median & $0.0000$ & \\
  quantiles $F^{-1}(u)$             & $0.0000$ & \\
  expectiles $e_\beta$              & $0.0000$ & \\
  \bottomrule
\end{tabular}
\end{center}

\textbf{The atom probabilities fail worst, and in the one direction this design cannot absorb.}
They are not a reading of the return at all --- a probability is not on the return axis --- and
on a fixed support a translate is a shift along the atom \emph{index}, so two well-separated
actions disagree maximally bin by bin. The reading therefore \emph{grows} the better the ordering
is resolved: $0$, $0.0087$, $0.0161$, $0.0246$, $0.0271$ as the separation goes $0$ to $8$. A
rule dividing by it would spread the policy widest exactly where one action is clearly best,
which inverts \S\ref{sec:rule}'s failure paragraph.

\paragraph{Scale, which is all that survives admissibility.} Two admissible families disagree
about $\srank$ by a factor, and whether that factor is constant decides whether $g$ absorbs it.
Uniform over levels against the two-point pair, by the shape of the centred advantage:

\begin{center}
\small
\begin{tabular}{lccc}
  \toprule
  advantage shape & over levels & two-point & ratio \\
  \midrule
  linear in $u$            & $0.2887$ & $0.2500$ & $1.15\times$ \\
  step at the median       & $1.0000$ & $1.0000$ & $1.00\times$ \\
  S-curve, heavier tails   & $0.8175$ & $0.7698$ & $1.06\times$ \\
  bottom $25\%$ only       & $0.4330$ & $0.2500$ & $1.73\times$ \\
  catastrophic bottom $5\%$& $0.2179$ & $0.0500$ & $\bm{4.36\times}$ \\
  jackpot top $5\%$        & $0.2179$ & $0.0500$ & $\bm{4.36\times}$ \\
  \bottomrule
\end{tabular}
\end{center}

\textbf{Two regimes, and only one is a problem.} Where the ordering disagreement is spread
broadly the families differ by $1.00$--$1.15\times$, a constant, absorbed entirely into $g$:
$\gapbar/\srank$ would move from $11.2$ to ${\sim}12.9$ and nothing else would. Where it is
concentrated in a tail they differ by $4.36\times$, and the factor is \emph{state-dependent}, so
it reshapes $\tau(s)$ rather than rescaling it and no choice of $g$ repairs it. Which regime a
game is in is a property of the game: a state where one action is fine except that it
occasionally ends the episode is the tail-concentrated case, which is \S\ref{sec:c2}'s. So the
coarser family is safe exactly where \S\ref{sec:c2} is not binding, and \S\ref{sec:owed} carries
the measurement that would say which.

Quantiles against expectiles on one risky-against-safe pair at equal means --- real disagreement,
both admissible --- reads $0.9522$ and $0.6098$, a factor of $0.64$. That is the correction of
\S\ref{sec:bg} priced: the family the head has been fitting all along differs from the one this
document described by a constant, and $g \approx 4$ was calibrated through it either way.

\paragraph{Quantisation, under a fixed support.} $\srank$ is small, so a reading that resolves no
finer than one atom spacing could be reporting nothing. Exact translates again, shifted by
fractions of the spacing, with C51's two-point projection applied as the head would:

\begin{center}
\small
\begin{tabular}{lccc}
  \toprule
  shift (atom spacings) & quantiles & expectiles & own median \\
  \midrule
  $0.10$ & $0.0000$ & $0.0006$ & $0.0005$ \\
  $0.25$ & $0.0000$ & $0.0012$ & $0.0014$ \\
  $0.50$ & $0.0000$ & $0.0021$ & $0.0027$ \\
  \bottomrule
\end{tabular}
\end{center}

Against a $\srank$ measured at $0.0171$--$0.0736$ (\S\ref{sec:denoms}) the spurious floor is
$8$--$120\times$ below the signal, so \textbf{quantisation is not a threat to the quantity}. But
the zeros in the first column are not a pass: the quantile reading is \emph{blind} below one
spacing rather than quiet, and on a game whose $\srank$ is small against the atom spacing it
would collapse to zero and turn \eqref{eq:rule} into argmax without saying so.
\eqref{eq:expectile} is smooth in $p$ and does not.

\paragraph{The knee is never reached, which is what picks the family.}
\texttt{IQN\_HUBER\_LOSS\_K} is $1$ against a raw $|td|$ that spans $0.02$--$0.45$ over a run, so
the quantile Huber sits in its quadratic branch throughout and \S\ref{sec:family}'s correction
holds for the whole of a budget rather than for part of it. That the estimand does not drift
between branches is the one thing here that a run could have spoiled and does not.

\paragraph{The closed form.} \eqref{eq:expectile} against a bisection on the same distributions,
maximum absolute error $4\times10^{-16}$ --- so the expectile family costs a
\texttt{searchsorted} and a division, not a root-find, which is what makes it affordable at the
same price as the quantile lookup it replaces.


\subsection{Action gaps}
\label{sec:meas-gap}

\paragraph{Action gaps.} \texttt{results/probe\_*.json} puts Phoenix's top-two gap inside BTR's
own $0.180$--$0.282$ band, so there is no action-gap deficit to explain. \S\ref{sec:denoms}'s
table supersedes the rest of this reading and sharpens it: the gap is a property of the game by
$5.7\times$ and of the behaviour policy by $1.3\times$, which is \S\ref{sec:project}'s second
constraint measured, and it is the movement that \eqref{eq:rule}'s denominator has to absorb.


\subsection{Exploration regret}
\label{sec:meas-rho}

\paragraph{Exploration regret.} \texttt{results/rho\_phoenix\_*.json}, $2048$ replay states per
checkpoint, at the per-decision horizon $H = 333$: the parameter draw costs $1.0$--$1.5\%$ of
$Q^\star$, the softmax at $\tau = 0.03$ costs $4.9$--$6.0\%$, and BTR's own $\epsilon = 0.01$
costs $11.5$--$17.5\%$ \emph{while winning}. Four checkpoints agree to some $30\%$, so the
quantity prices the policy rather than the weights. Below $\tau \approx 0.003$ the series flattens
onto the draw's floor --- which is \S\ref{sec:meas-g}'s point, and it is also why no controller can
usefully servo $\rho$ on this codebase without a second value head to price the draw against.


\subsection{Return: full runs, and $J(\tau)$ at fixed weights}
\label{sec:meas-return}

\paragraph{Four full runs.} $200$M-frame Phoenix arms differing in acting sharpness, one seed
each, scored against BTR's \texttt{results.csv} Phoenix row, geometric mean:

\begin{center}
\begin{tabular}{lcccc}
  \toprule
  arm & $1$--$50$M & $50$--$200$M & $100$--$200$M & $150$--$200$M \\
  \midrule
  argmax + parameter draw & $\bm{0.693}$ & $\bm{1.072}$ & $\bm{1.057}$ & $\bm{1.128}$ \\
  softmax, $\tau = 0.005$ & $0.304$ & $0.901$ & $0.945$ & $0.968$ \\
  softmax, $\tau = 0.01$  & $0.213$ & $0.718$ & $0.818$ & $0.935$ \\
  softmax, $\tau = 0.03$  & $0.630$ & $0.787$ & $0.825$ & $0.910$ \\
  \bottomrule
\end{tabular}
\end{center}

Three readings, and they set every priority above. \emph{Each soft arm is the argmax curve
delayed}, by $22.5$M, $42.5$M and $40$M frames respectively, overlaying to within $1$--$6\%$ after
the shift --- so sharpness buys and spends time rather than choosing a point on a trade-off curve.
\emph{The level is worth less than a seed}: $19\%$ spread over $150$--$200$M against $3.2\times$
over $1$--$50$M, with one seed per arm, which is why no controller on the level appears in this
document. \emph{All of the damage is early}, which is where \S\ref{sec:c5} and \S\ref{sec:c6}
live and where nothing else on the project's plan is aimed.

\paragraph{$J(\tau)$ at fixed weights.} $J(\tau)$ is the return of the policy at acting scale
$\tau$ played on weights that are not being updated --- what a level criterion would be servoing,
and the only thing here measured by rollout. \texttt{scripts/score\_policies.py}, two Phoenix
checkpoints, $32$ episodes per cell, returns in thousands with the standard error of the mean:

\begin{center}
\small
\setlength{\tabcolsep}{4.5pt}
\begin{tabular}{lccccc}
  \toprule
  weights & argmax & $0.003$ & $0.01$ & $0.03$ & $0.1$ \\
  \midrule
  soft-trained, $\tau = 0.03$ & $\bm{368.8}\pm27.8$ & $340.5\pm27.3$ & $310.0\pm34.0$
                              & $300.7\pm31.7$ & $164.5\pm16.4$ \\
  argmax-trained              & $200.0\pm16.6$ & $169.8\pm17.1$ & $\bm{227.2}\pm23.0$
                              & $169.9\pm19.7$ & $56.6\pm5.6$ \\
  \bottomrule
\end{tabular}
\end{center}

The signal is $2$--$5\%$ per octave and the measurement noise is $8$--$12\%$, which is the reason
no criterion in this document scores a policy by rollout. On the soft-trained weights $J$ falls
monotonically and argmax wins, which is why \eqref{eq:rule} is built to contain argmax rather than
to improve on it. The cells do not carry a normal approximation --- the soft-trained $0.01$ cell
runs from $7{,}240$ to $583{,}990$ --- and the shape is suggestive at best.


\subsection{The off-greedy rate and the gain $g$}
\label{sec:meas-g}

The rule's two constants are dimensionless and each has an anchor rather than a range to
search. This subsection is the first of them.

\paragraph{$g$, and what it is measured to buy.} One gain, four checkpoints, off-greedy
probability under \eqref{eq:rule} against the fixed temperature it replaces:

\begin{center}
\small
\begin{tabular}{lccccc}
  \toprule
  & Phoenix $46$M & Phoenix $250$M & Phoenix $200$M & Zaxxon $204$M & spread \\
  \midrule
  fixed $\tau = 0.03$      & $0.169$ & $0.095$ & $0.090$ & $0.476$ & $5.3\times$ \\
  \eqref{eq:rule}, $g = 4$ & $0.049$ & $0.053$ & $0.044$ & $0.057$ & $1.3\times$ \\
  \eqref{eq:rule}, $g = 6$ & $0.032$ & $0.037$ & $0.029$ & $0.037$ & $1.3\times$ \\
  \bottomrule
\end{tabular}
\end{center}

\emph{This is the whole claim of the document, in one table.} A constant in value units delivers
a $5.3\times$ range of behaviour across two games, two behaviour policies and two points in
training; the same rule at one dimensionless gain delivers a $1.3\times$ range over the same set.
The relative regret $\rho/\gapbar$ agrees to the same tolerance, $0.0016$--$0.0032$ at $g = 4$.

\textbf{$g \approx 4$}, then, from the one external anchor that survives: it is where the
undirected layer contributes about what the parameter draw already contributes on its own
($3.8$--$4.5\%$ off-greedy). $g \approx 3$ doubles the total spend, which is a defensible first
choice and a decision to log rather than a search. One caveat on the anchor, from
\S\ref{sec:related}: BTR's Table~5 measures $3.8\%$ policy churn on Phoenix, so the rate being
matched is the same size as the involuntary change in the greedy action between two updates, and
some of what the anchor prices may not be attributable to the behaviour policy at all. Note what this anchor is \emph{not} --- a
claim that $4.5\%$ is the right rate. The two best arms this project has run sit at $4\%$
(Phoenix argmax) and $48\%$ (Zaxxon at $\tau = 0.03$), so the rate that wins is not portable even
where the rule that produces it is. \eqref{eq:rule} makes the two games commensurable; it does
not say where on the axis to stand, and \S\ref{sec:readmit} is the only thing here that could.

\emph{Why not calibrate against the best soft arm's measured $\rho$.} Because $83$--$90\%$ of that
number is the parameter draw. On \texttt{phoenix\_2026-09-20\_02-21-56}'s final checkpoint the
draw alone reads $\rho = 7.35\times10^{-4}$ and the $\tau = 0.005$ policy reads
$8.9\times10^{-4}$, so the softmax's own contribution is $1.6\times10^{-4}$ --- a sixth of what a
naive match would target. The previous revision's staging instruction matched the two directly
and would have overspent by that factor. Net the floor out, always.


\subsection{Sharpness and the floor $k_0$}
\label{sec:meas-k}

The floor is stated in the sharpness of \eqref{eq:rule-k}, so that $k_0$ is directly comparable
with the arms that have run and with the fixed-weight cells of \S\ref{sec:meas-return}. That
comparison is what moved it by an order of magnitude.

\paragraph{The previous anchor was read off a model, and the model does not hold.} $k_0 = 4.4$
came from a state with one best action and the rest at a common gap. Real states are not that:
the top-two gap is $0.187$ against $\gapbar = 0.721$ on Phoenix, so a near-tie and not the mean
gap sets the off-greedy rate, and the tail of clearly bad actions inflates $\gapbar$ without
contributing any probability. Scored on real states instead of on the model, $k = 4.4$ puts some
$40\%$ of decisions off the greedy action --- more diffuse than $\tau = 0.1$, the fixed-weight
cell that scored $164$k against argmax's $369$k. The arms that actually ran sit at $k = 24$
($\tau = 0.03$) to $k = 144$ ($\tau = 0.005$) on Phoenix. \textbf{The working band is $k = 24$ to
$300$, and the previous revision's floor was an order of magnitude below the bottom of it}: it
would not have guarded the rule, it would have been the rule.

The same scoring validates the method, which is worth one line: at $\tau = 0.03$ it gives an
off-greedy rate of $0.095$ against the $0.102$--$0.104$ measured directly on the acting stream.

\paragraph{$k_0$, as a guard and not an operating point.} The floor binds where
$\gapbar/\srank < k_0/g$, and that ratio's $\text{p}5$ is $4.5$ on Phoenix and $4.7$ on Zaxxon.
So $k_0 \approx 20$ --- the sharpness of the most diffuse arm that has ever survived a full
budget --- rests a few per cent of states on the floor at $g = 4$ and fewer at $g = 6$, which is
what a guard should do. Log the resting fraction; if it climbs, $\srank$ is failing and
\S\ref{sec:rule}'s failure paragraph is being exercised.

\emph{Measured, the floor behaves. The ceiling behaves at one value of a constant and not at
another, which is the finding.} At $g = 4$ and the $\tau_\ell = 0.03$ in force, the floor takes
$3.2$--$6.0\%$ of states across the four checkpoints and the ceiling $\tau(s) \le \tau_\ell$ takes
$0.1\%$, $3.8\%$, $3.7\%$ and $0.4\%$: both are guards doing their job. Lower $\tau_\ell$ to the
$0.02$ that one run used and the ceiling's share becomes $1.3\%$, $26.7\%$, $17.4\%$ and $0.8\%$,
so on the two later Phoenix checkpoints it stops being a guard and becomes the operating branch
for a sixth to a quarter of decisions --- over which \eqref{eq:rule} is the fixed temperature it
exists to replace.

\textbf{A $1.5\times$ move in a loss constant changes the rule's resting fraction by $7\times$},
and that is a property of the decoupled design rather than of either value: the ceiling is the one
place the behaviour policy touches the loss (\S\ref{sec:bg}), so the rule inherits whatever is
done to $\tau_\ell$ for reasons of its own. \S\ref{sec:decouple-ledger} carries it as decoupling's
cost and \S\ref{sec:couple-case} as the reason coupling deletes the branch outright. $g = 6$
holds the ceiling under $5\%$ at both values --- $0.2\%$, $4.3\%$, $4.1\%$, $0.6\%$ at $0.02$ ---
which is one more reason to prefer the higher gain and is the only place in this document where
the two gains are separated by something other than taste.

\subsection{The distortion a coupled $\tau$ induces}
\label{sec:meas-couple}

\S\ref{sec:c8} weighs putting $\tau(s)$ into the loss, and \eqref{eq:distortion} is the whole of
what that changes: a per-step intrinsic reward $D(s)$, in the units of $\tilde{q}$.
\texttt{scripts/coupling\_cost.py} evaluates it on the same four checkpoints' per-state arrays
and writes \texttt{results/coupling\_cost.json}; it is a post-process on
\texttt{spread\_probe.py}'s output, so it needs no checkpoint and re-runs in a second. Four
rows per checkpoint: \eqref{eq:rule} at $g = 4$; the same states held at the fixed
$\tau_\ell = 0.03$, which is both the ceiling and what every arm in this section actually ran at,
so the two coincide and only one row is needed; and candidate (D), $\tau = \gapbar/k_0$ at the
$k_0 = 4.4$ the previous revision proposed, which is the rule \eqref{eq:rule} was ranked
\emph{below} on this question.

\begin{center}
\small
\setlength{\tabcolsep}{4.5pt}
\begin{tabular}{llcccccccc}
  \toprule
  & & & & & \multicolumn{2}{c}{$D/\gapbar$} & \multicolumn{3}{c}{lift $\E_D[x]/\E[x]$} \\
  \cmidrule(lr){6-7}\cmidrule(lr){8-10}
  & & $k$ & $D$ & $D/Q^\star$ & mean & p95 & top-two & $\srank$ & $\gapbar$ \\
  \midrule
  Phoenix $46$M  & \eqref{eq:rule}, $g=4$ & $47$  & $5.6{\cdot}10^{-4}$ & $3.1{\cdot}10^{-5}$ & $0.19\%$ & $1.3\%$  & $\bm{0.11}$ & $1.20$ & $0.86$ \\
                 & $\tau = \tau_\ell$     & $13$  & $6.4{\cdot}10^{-3}$ & ---                 & $2.4\%$  & $10.0\%$ & $0.25$ & $0.92$ & $0.74$ \\
                 & (D), $k_0 = 4.4$       & $4.4$ & $4.0{\cdot}10^{-2}$ & $2.2{\cdot}10^{-3}$ & $\bm{11.0\%}$ & $20.0\%$ & $0.82$ & $1.06$ & $\bm{1.07}$ \\
  \midrule
  Phoenix $250$M & \eqref{eq:rule}, $g=4$ & $46$  & $9.5{\cdot}10^{-4}$ & $4.4{\cdot}10^{-5}$ & $0.16\%$ & $1.0\%$  & $\bm{0.09}$ & $1.11$ & $0.89$ \\
                 & $\tau = \tau_\ell$     & $24$  & $3.5{\cdot}10^{-3}$ & ---                 & $0.58\%$ & $2.8\%$  & $0.15$ & $0.86$ & $0.88$ \\
                 & (D), $k_0 = 4.4$       & $4.4$ & $7.7{\cdot}10^{-2}$ & $3.3{\cdot}10^{-3}$ & $\bm{11.1\%}$ & $20.5\%$ & $0.70$ & $0.97$ & $\bm{1.01}$ \\
  \midrule
  Phoenix argmax & \eqref{eq:rule}, $g=4$ & $67$  & $7.2{\cdot}10^{-4}$ & $3.4{\cdot}10^{-5}$ & $0.11\%$ & $0.69\%$ & $\bm{0.09}$ & $1.15$ & $0.80$ \\
                 & $\tau = \tau_\ell$     & $31$  & $3.3{\cdot}10^{-3}$ & ---                 & $0.54\%$ & $2.5\%$  & $0.14$ & $0.85$ & $0.78$ \\
                 & (D), $k_0 = 4.4$       & $4.4$ & $1.1{\cdot}10^{-1}$ & $4.7{\cdot}10^{-3}$ & $\bm{12.4\%}$ & $21.1\%$ & $0.82$ & $0.93$ & $\bm{1.07}$ \\
  \midrule
  Zaxxon $204$M  & \eqref{eq:rule}, $g=4$ & $45$  & $3.3{\cdot}10^{-4}$ & $3.0{\cdot}10^{-5}$ & $0.23\%$ & $1.4\%$  & $\bm{0.33}$ & $3.15$ & $1.50$ \\
                 & $\tau = \tau_\ell$     & $5.3$ & $2.2{\cdot}10^{-2}$ & ---                 & $22.1\%$ & $72.5\%$ & $0.60$ & $0.83$ & $0.78$ \\
                 & (D), $k_0 = 4.4$       & $4.4$ & $2.6{\cdot}10^{-2}$ & $2.3{\cdot}10^{-3}$ & $\bm{16.3\%}$ & $24.8\%$ & $1.01$ & $1.61$ & $\bm{1.34}$ \\
  \bottomrule
\end{tabular}
\end{center}

\paragraph{The shape of $D$, which is what \S\ref{sec:bg} reads it by.} Same states, same script,
under \eqref{eq:rule} at $g = 4$. The identity $D = \tau\h(\pi_\tau) - \rho(\pi_\tau)$ holds to
$10^{-14}$ on all four, and $\rho/\tau\h$ averages $0.57$--$0.83$, so the free bound
$\rho \le \tau\h \le \tau\ln|\actions|$ is real but loose: against $\tau\ln|\actions|$ the realised
regret is $1.6\times10^{-4}$ to $1.3\times10^{-3}$ at the median state and $0.16$--$0.19$ at the
$99$th percentile. $D/\tau$ has a median of $4.3\times10^{-5}$ to $4.8\times10^{-4}$ and a mean
$132$ to $1137\times$ larger, which is the tail this subsection's weighted means are reading.

\paragraph{How many actions are in contention, and what the top-two shortcut costs.} The
contention count $m = e^{D/\tau}$ of \S\ref{sec:bg} has a median of $1.0000$--$1.0005$ --- the
typical state has one action in play and earns nothing --- but a $D$-weighted mean of
$1.49$--$1.56$, and $46$--$53\%$ of $\sum D$ sits on states with $m > 1.5$. \textbf{So wherever
the bonus is actually paid, the runner-up is not alone}, and the top-two form
$D \approx \tau(1 + e^{-k_2})$ is being used exactly where it is weakest. Scored properly it
costs: a median relative error of $0.00$--$1.48\%$, which is the figure a careless reading would
quote, against a \textbf{$D$-weighted error of $3.3$, $4.3$, $3.3$ and $13.0\%$} and a $99$th
percentile of $44$--$51\%$. The median is a statistic of the resolved states, where both forms
agree at nothing. The larger action set is the worse case, as the mechanism predicts. For
reference the near-tie sharpness $k_2 = \gap_2/\tau$ reads $0.66$--$0.84$, $7.7$--$10.2$ and
$24$--$37$ at the fifth, fiftieth and ninety-fifth percentiles.

\paragraph{What duplicating the action set would add.} $D = \tau\ln m$ and $m$ counts actions, so
a duplication sends $D \mapsto D + \tau\ln(\text{copies})$ exactly --- confirmed here to
$4\times10^{-16}$ relative, at two and three copies, on both games. The size is the point:
\textbf{one duplication adds $8.4\times$ the entire mean $D$ on Zaxxon and $11.1\times$ it on
Phoenix.} The added term is $\tau\ln 2$, and on the operating branch $\tau = \srank/g$, so it
correlates with $\srank$ at $0.78$--$0.92$ across these states --- a bonus proportional to the
denominator, arriving from the action set's spelling rather than from any near-tie.
\S\ref{sec:c8} is what that does to the character of a coupled rule, and \S\ref{sec:c3} is the
case the behaviour rule passes and the coupled objective does not.

\paragraph{Two constants this subsection is read against.} \eqref{eq:couplegain}'s sufficient
gain, evaluated at $\gamma = 0.997$ and $\alpha = 0.9$, is $g \gtrsim 69L$ on Phoenix and
$96L$ on Zaxxon, with $L \le 2$; against the $g \approx 4$ of \S\ref{sec:meas-g} the cheap
contraction argument misses by a factor of $17$ to $48$, and \S\ref{sec:couple-fixedpoint} is
what follows. And \texttt{IQN\_ACT\_SAMPLES} is $32$ against \texttt{IQN\_TRAIN\_SAMPLES}'s $8$,
so a coupled $\hat\tau$ read in the loss carries about twice the relative sampling error of the
acting one, $27\%$ against $13\%$.

\paragraph{Which branch pays it, what one ratio decides it, and how hard it pulls.}
\eqref{eq:rule} is a $\min$ of three, and ``$\tau$ proportional to $\srank$'' describes only the
first, so the share of $\sum D$ each branch carries is what says whether the description covers
the measurement. It does: the $\srank/g$ branch carries $87\%$, $89\%$, $84\%$ and $67\%$, the
floor $13\%$, $7\%$, $7\%$ and $24\%$, the ceiling $0\%$, $4\%$, $9\%$ and $9\%$. On that branch
$D \approx (\srank/g)\exp(-g\,\gap_2/\srank)$, so the controlling quantity is the ratio
$\gap_2/\srank$ and not either part of it --- and the lifts agree, $0.09$, $0.09$, $0.08$ and
$0.11$ for the ratio against $0.09$--$0.33$ for $\gap_2$ and $1.11$--$3.15$ for $\srank$.
$\sret$, the return width the risk-seeking charge was really about, lifts $1.13$, $1.07$, $1.08$
and $1.53$: mild on Phoenix, real on Zaxxon.

\emph{As a standing bias rather than a per-step term}, at the per-decision horizon of
\S\ref{sec:meas-rho}, $H\,D/Q^\star$ is $1.05\%$, $1.36\%$, $1.04\%$ and $0.95\%$ --- the same
bracket as the parameter draw's own $1.0$--$1.5\%$ exploration regret. The per-step comparison
below and this one answer different questions and both belong in \S\ref{sec:c8}.

\paragraph{Magnitude.} \textbf{Under \eqref{eq:rule} the induced intrinsic reward is $0.11$ to
$0.23\%$ of the mean action gap and $3$--$4\times10^{-5}$ of $Q^\star$.} Against that,
the cost of \emph{not} coupling --- the Munchausen bonus evaluated at an action the loss did not
anticipate, which \S\ref{sec:decouple-ledger} accepts --- reads $-0.043$ against $-0.007$ across a
twofold range of behaviour entropy on Zaxxon, or about $0.46\%$ of that run's
$\texttt{train/q} \approx 7.8$. The two are per-step additive terms in the same units, so the
comparison is direct and it goes the wrong way for \S\ref{sec:c8}'s original argument. Zaxxon
carries both: the lost cancellation is $0.036$ in value units and $D$ under the rule is
$3.3\times10^{-4}$, so coupling distorts the objective by \emph{one part in $109$} of what
decoupling costs --- $0.0042\%$ of the critic's own scale against $0.46\%$. Both figures are at
the $\tau_\ell$ in force, so neither needs rescaling to sit beside the other. What argues against coupling is
\S\ref{sec:couple-fixedpoint} and \S\ref{sec:c8-staging}, neither of which depends on a
magnitude. The intrinsic reward is not the reason and, at this size, is an argument the other
way.

\paragraph{The magnitude is set by sharpness, so the ranking of the candidates inverts.} Read
down any block: $D/\gapbar$ rises $4$--$13\times$ on Phoenix and $98\times$ on Zaxxon as $k$
falls from the rule's $45$--$67$ to the fixed temperature's $5.3$--$31$. $D$ is linear in $\tau$
and exponential in $-\gap/\tau$ (\S\ref{sec:bg}), and the exponential is what moves.
\eqref{eq:rule} is the sharpest rule in the family, so it is the one whose coupled bonus is
negligible; candidate (D) at $k_0 = 4.4$, the diffuse end of the same family, averages
$11$--$16\%$ of the mean gap with a $\text{p}95$ of $20$--$25\%$, and the fixed $\tau = \tau_\ell$
on Zaxxon reaches $72\%$ at $\text{p}95$. The diffuse rules are where the objection has force, and
every one of them is a rule this document has already retired for other reasons.

\paragraph{Carrier: the bonus is paid on the near-tie, not on the denominator.} The last three
columns weight each per-state quantity by the $D$ it carries and divide by its plain mean ---
$\E_D[x]/\E[x]$, a lift of $1$ meaning no concentration. That is a weighted mean rather than a
correlation of $\ln D$ deliberately, and the reason is sharper than ``$D$ is small''.
$D = \tau\ln(1 + e^{-k_2} + \cdots)$ is a $\ln(1+x)$ whose $x$ falls below float64's resolution as
soon as $k_2 > 36.7$, so a direct evaluation returns \emph{exactly} zero --- on $4.4\%$, $2.6\%$,
$5.8\%$ and $0.3\%$ of these states --- where the true value is of order $10^{-51}\tau$. $\ln D$ is
then $-\infty$ on up to a twentieth of the sample, and two algebraically identical evaluations
agreeing to $5\times10^{-18}$ give $-0.35$ and $-0.70$ for the same log-correlation. The weighted
mean is immune because those states carry no weight, and that is checkable rather than hoped for:
every figure in the table above is identical to sixteen digits under a \texttt{log1p} evaluation.

Under \eqref{eq:rule} the mass sits on states whose top-two gap is a tenth of average on Phoenix
and a third on Zaxxon, and it is concentrated: the top decile of states carries $73.5$, $76.5$,
$80.5$ and $79.5\%$ of $\sum D$. $\srank$, which the rule divides by and which \S\ref{sec:c8}
originally named as what the agent would be paid to seek, shows next to no concentration on
the three Phoenix checkpoints ($1.11$--$1.20$) and a real one ($3.15$) only on Zaxxon. So the induced reward pays
the agent to stand where the two best actions are nearly tied in mean --- the cheapest place in
the state space to be paid to stand, the regret of a wrong choice there being small by
construction --- and the risk-seeking reading does not survive the measurement on either game.

\emph{Candidate (D)'s row is the control, and it behaves exactly as the previous revision said
(F) would.} At $k_0 = 4.4$ the bonus is proportional to the denominator --- $\gapbar$'s lift is
$1.01$--$1.34$, against \eqref{eq:rule}'s $0.80$--$1.50$ --- it is spread rather than
concentrated, its top decile carrying only $18$--$25\%$, and at $11$--$16\%$ of $\gapbar$ it is
$50$--$110\times$ larger. \textbf{So the objection is sound, it was written about the right
mechanism, and it was attached to the wrong candidate.} A rule that holds $k$ constant holds
$\h(\pi)$ roughly constant with it, which is what makes its bonus track its denominator; a rule
whose $k$ varies, as \eqref{eq:rule}'s does over the measured range of $\gapbar/\srank$, hands the
exponential the decision instead.

\paragraph{The Munchausen clip, counted two ways.} $l_0$ fires at $(s,a)$ iff
$\gap_a + D > 1/\alpha$, so counted uniformly over state-action pairs at $\tau_\ell$ it fires on
$3.5\%$, $22.0\%$ and $42.7\%$ of the Phoenix pairs and $0.1\%$ of the Zaxxon ones --- not the
``never'' the previous revision recorded. Weighted by the behaviour that draws them, which is what
the loss sees, it is zero to underflow under \eqref{eq:rule} and $3\times10^{-4}$ to
$5\times10^{-3}$ even for candidate (D), whose $\tau$ reaches $0.09$--$0.21$. The previous
revision's projection of $19\%$ of state-action pairs for its own coupled rule was the uniform
count, which is not the distribution the loss draws from. \S\ref{sec:couple-case} is the identity
that makes the second number small.

\paragraph{Gating, and it is the finding with consequences beyond \S\ref{sec:c8}.} $D \le
\tau\ln|\actions|$, and how much of that ceiling a state receives is set by its own gaps. Hold
\emph{every} state at $\tau_\ell$ --- the most a coupled rule is ever permitted to pay --- and the
fraction still receiving under a tenth of $\tau_\ell\ln|\actions|$ is $62.5\%$, $79.3\%$ and
$80.1\%$ on Phoenix and $15.3\%$ on Zaxxon; under the rule itself it is $95$--$99\%$ on Phoenix.
\textbf{The entropy channel delivers nothing to a state whose gaps are large, however unresolved
that state actually is.} A reward-level bonus --- RND's, a count's --- is ungated: it is added
whatever the critic believes. This bounds what \S\ref{sec:coupling} can buy before any run is
spent, and it is the one reading here that is about family (I) rather than about
\eqref{eq:rule}.

\paragraph{What this measurement is not.} Every figure is taken at the gaps a \emph{decoupled}
net learned, so it is a linearisation of a scheme that would move them --- Remark~\ref{rem:coupled-ok}
is why that caveat is load-bearing rather than formal. It prices the objective's distortion and
says nothing about the optimisation: a term two orders of magnitude below the gaps can still
matter if it is systematic, and the only thing that settles that is the arm.


\subsection{Staging}
\label{sec:staging}

\paragraph{Staging.} Three steps, the first two free.

\begin{enumerate}
  \item \textbf{Log the statistics in a training run} --- $\srank$, $\gapbar$, the percentiles
  of their ratio and the floor's resting fraction --- under the unchanged behaviour policy. Free,
  on the control path only as a diagnostic, and it answers the one question
  \S\ref{sec:denoms} could not: the four checkpoints are four runs, and whether
  $\gapbar/\srank$ is stable \emph{within} one run has not been seen.
  \item \textbf{Score \eqref{eq:rule} at fixed weights} with \texttt{scripts/score\_policies.py},
  as a cell beside the existing argmax / $0.003$ / $0.01$ / $0.03$ columns, at two or three
  $(g, k_0)$. This is the same screen that retired the level criteria, applied to the rule that
  replaced them. Its limit must be stated with its result: a fixed-weight readout says nothing
  about which behaviour \emph{produces} better weights, so it is a screen and not a verdict. The
  bar is to be within noise of the argmax cell --- $8$--$12\%$ standard error at $32$ episodes ---
  not to beat it.
  \item \textbf{One run}, against \texttt{phoenix\_2026-09-20\_02-21-56} so the comparison is
  soft-against-soft at the best soft arm. Read it on $1$--$50$M, which is where the arms differ by
  $3.2\times$ and where every case above says the rule binds. Log $k(s)$'s percentiles, the
  fraction of states on the floor branch, the fraction on the $\tau_\ell$ ceiling, $\srank$'s
  trajectory, and $\rho$ against the draw-only floor. Those five say which branch is operating and
  whether either constant is misplaced, which nothing visible from the acting policy alone can.
\end{enumerate}

\section{Open questions}
\label{sec:owed}

\begin{itemize}
  \item \textbf{\S\ref{sec:fixdraw}, before anything else.} Every epistemic denominator
  measured --- $\dbar$, \eqref{eq:uasplit}'s $\tilde\sigma_{\mathrm{epi}}$, and Fox's squared
  form built on it --- fails to transfer, and \S\ref{sec:denoms} locates the failure in the
  draws rather than in the statistic. Nothing about \S\ref{sec:c10} or \S\ref{sec:c11} can be
  decided until the draws are posterior samples, and after that all three are re-measurable
  offline in minutes.
  \item \textbf{Fox's exponent, once there is an epistemic $\sigma$ worth squaring.} Linear
  against squared in the signal-to-noise ratio is a larger difference than anything else in this
  document, and it is untestable while the denominator is the problem.
  \item \textbf{Is the ordering disagreement tail-concentrated?} \S\ref{sec:meas-family} finds
  the choice of family costing a constant factor when the centred advantage varies broadly and a
  state-dependent one when it lives in a tail, and only the first is absorbed by $g$. The
  measurement is
  the distribution of $A_\beta(s,a)$ along the level axis on the existing checkpoints --- already
  on disk as \texttt{results/spread\_*.npz}'s inputs --- and it decides whether a categorical head
  may use a cheap two-point family or must carry a fuller one. It is a prerequisite for the
  head, not a question about the temperature.
  \item \textbf{$g$ across the change of head.} Admissibility is settled (\S\ref{sec:meas-family})
  but scale is not (\S\ref{sec:meas-family}), so $g$ has to be re-measured whenever the family
  moves --- including if
  \texttt{IQN\_HUBER\_LOSS\_K} is lowered to recover the quantile estimand, which is an open item
  in \texttt{btr.py} that was never connected to this one.
  \item \textbf{$\gapbar/\srank$ inside a single run.} The four checkpoints of
  \S\ref{sec:denoms} are four runs. That the ratio is stable across them is this document's
  central empirical claim and it has never been watched as one run moves. \S\ref{sec:staging}'s
  first stage is this and it is free.
  \item \textbf{A third game.} Two games agreeing closely (\S\ref{sec:denoms}) is two games, and
  this is the claim
  most worth trying to break.
  \item \textbf{Whether $\sret$ is a return dispersion at all.} It is flat across Phoenix
  checkpoints while \texttt{quantile\_embedding}'s bias grows by orders (\S\ref{sec:denoms}).
  Compare it
  against the realised dispersion of Monte-Carlo returns from the same states.
  \item \textbf{\eqref{eq:rule} scored at fixed weights} beside the existing $J(\tau)$ cells ---
  \S\ref{sec:staging}, stage 2 --- and \textbf{$J(\tau)$ on a soft-trained Zaxxon net}, both curves
  so far being Phoenix.
  \item \textbf{\eqref{eq:distortion} evaluated on a contracting $\hat{s}$}, before any coupled
  arm is run. \S\ref{sec:meas-couple} shows the entropy channel pays almost nothing to a state
  with large gaps, so the question of whether a novelty signal survives that gating is answerable
  in minutes with \texttt{scripts/coupling\_cost.py} and decides whether \S\ref{sec:coupling} is
  worth a baseline at all. Behind \S\ref{sec:fixdraw}, like everything else epistemic.
  \item \textbf{Whether a $\lambda$ read off the iterate converges at all}, which
  \S\ref{sec:couple-fixedpoint} makes the gating item for every coupled scheme here and which
  nothing in this document answers. Two jobs, both small, either informative alone:
  recompute $\lambda(s)$ from the iterate on every sweep of \texttt{scripts/check\_mvi\_gap.py}'s
  random finite MDPs and see whether the composed map converges and how that depends on $g$; and
  measure $\srank$'s Lipschitz constant $L$ over replay states, which turns
  \eqref{eq:couplebound} from indicative into quantitative. Everything about coupling waits on
  the first of these and it is the cheapest item on this list.
  \item \textbf{A coupled arm on a contracting denominator} (\S\ref{sec:coupling}), if the
  previous two items clear. It is the only thing in this document that would make the sampling
  rule into directed exploration rather than better-shaped dithering, and it is the one item here
  that costs a baseline.
  \item \textbf{The ceiling's sensitivity to $\tau_\ell$.} \S\ref{sec:meas-k} finds the ceiling
  turning from a guard into a major operating branch under a modest change in $\tau_\ell$, a
  constant chosen for the loss's reasons. Nothing anticipated the interaction and no run has
  logged it. A higher gain holds it to a guard at either value and is the cheapest fix; coupling
  removes the branch outright, which \S\ref{sec:couple-case} counts as one of its four arguments.
  \item \textbf{Boltzmann--Gumbel in a discounted control problem.} \S\ref{sec:related} confirms
  the bound is a stochastic-bandit result. Whether the $\sqrt{N_a}$ law or the suboptimality of
  monotone schedules survives bootstrapping is not addressed anywhere read here, and it is the
  only formal result near this design.
\end{itemize}

\end{document}
